% Options for packages loaded elsewhere
\PassOptionsToPackage{unicode}{hyperref}
\PassOptionsToPackage{hyphens}{url}
\documentclass[
  11pt,
]{article}
\usepackage{xcolor}
\usepackage[margin=1in]{geometry}
\usepackage{amsmath,amssymb}
\setcounter{secnumdepth}{-\maxdimen} % remove section numbering
\usepackage{iftex}
\ifPDFTeX
  \usepackage[T1]{fontenc}
  \usepackage[utf8]{inputenc}
  \usepackage{textcomp} % provide euro and other symbols
\else % if luatex or xetex
  \usepackage{unicode-math} % this also loads fontspec
  \defaultfontfeatures{Scale=MatchLowercase}
  \defaultfontfeatures[\rmfamily]{Ligatures=TeX,Scale=1}
\fi
\usepackage{lmodern}
\ifPDFTeX\else
  % xetex/luatex font selection
  \setmainfont[]{Cambria}
  \setmonofont[]{Consolas}
  \setmathfont[]{Cambria Math}
\fi
% Use upquote if available, for straight quotes in verbatim environments
\IfFileExists{upquote.sty}{\usepackage{upquote}}{}
\IfFileExists{microtype.sty}{% use microtype if available
  \usepackage[]{microtype}
  \UseMicrotypeSet[protrusion]{basicmath} % disable protrusion for tt fonts
}{}
\makeatletter
\@ifundefined{KOMAClassName}{% if non-KOMA class
  \IfFileExists{parskip.sty}{%
    \usepackage{parskip}
  }{% else
    \setlength{\parindent}{0pt}
    \setlength{\parskip}{6pt plus 2pt minus 1pt}}
}{% if KOMA class
  \KOMAoptions{parskip=half}}
\makeatother
\usepackage{longtable,booktabs,array}
\newcounter{none} % for unnumbered tables
\usepackage{calc} % for calculating minipage widths
% Correct order of tables after \paragraph or \subparagraph
\usepackage{etoolbox}
\makeatletter
\patchcmd\longtable{\par}{\if@noskipsec\mbox{}\fi\par}{}{}
\makeatother
% Allow footnotes in longtable head/foot
\IfFileExists{footnotehyper.sty}{\usepackage{footnotehyper}}{\usepackage{footnote}}
\makesavenoteenv{longtable}
\setlength{\emergencystretch}{3em} % prevent overfull lines
\providecommand{\tightlist}{%
  \setlength{\itemsep}{0pt}\setlength{\parskip}{0pt}}
% Center all longtables
\setlength{\LTleft}{0pt plus 1fill}
\setlength{\LTright}{0pt plus 1fill}
\usepackage{bookmark}
\IfFileExists{xurl.sty}{\usepackage{xurl}}{} % add URL line breaks if available
\urlstyle{same}
\hypersetup{
  hidelinks,
  pdfcreator={LaTeX via pandoc}}

\author{}
\date{}

\begin{document}

\begin{center}
{\LARGE\bfseries Character-Theoretic Effective Rank:\\[0.4em]
From Schur's Lemma to a Stereoscopic Measurement\\[0.2em]
Apparatus on Continuous Hamiltonian Substrates}

\vspace{1.5em}

{\large Antonio P.\ Matos}\\[0.4em]
{\small ORCID: \href{https://orcid.org/0009-0002-0722-3752}{0009-0002-0722-3752}}\\[0.3em]
April 24, 2026\\[0.3em]
{\small DOI: \href{https://doi.org/10.5281/zenodo.19744573}{10.5281/zenodo.19744573}}\\[0.3em]
Independent Researcher; Fancyland LLC / Lattice OS\\[0.3em]
Preprint v1.0

\vspace{0.8em}

{\small\textbf{MSC 2020:} 60J22, 20C30, 15A18, 62H25, 15B52, 82C22}

\vspace{0.4em}

{\small\textbf{Keywords:} Schur's lemma, effective rank, equivariant Markov generator, detailed balance, Wishart-Nyquist sampling, character-entropy decomposition, triangle-wave overtone theorem, holographic dictionary, Hamiltonian substrate, A-optimality, Heisenberg duality, manifold curvature, semantic interferometer}
\end{center}

\textbf{Companion papers:}

\begin{itemize}
\tightlist
\item
  {[}1{]} A. P. Matos, ``The Unity Clock: Effective Dimensional Collapse
  of the Addition-Multiplication Commutator on Coprime Lattices,''
  preprint (2026). DOI:
  \href{https://doi.org/10.5281/zenodo.19478727}{10.5281/zenodo.19478727}
\item
  {[}2{]} A. P. Matos, ``The Arithmetic Black Hole: Softmax
  Thermodynamics and the Four Eigenvalue Laws of the Prime Gas,''
  preprint (2026). DOI:
  \href{https://doi.org/10.5281/zenodo.19442006}{10.5281/zenodo.19442006}
\item
  {[}3{]} A. P. Matos, ``Active Transport on the Prime Gas: Flat-Band
  Condensation, the Rabi Phase Transition, and the Arithmetic Qubit,''
  preprint (2026). DOI:
  \href{https://doi.org/10.5281/zenodo.19243258}{10.5281/zenodo.19243258}
\item
  {[}4{]} A. P. Matos, ``Universal Two-Prime Formula for the
  Coprime-Lattice Coupling Constant,'' preprint (2026). DOI:
  \href{https://doi.org/10.5281/zenodo.19210625}{10.5281/zenodo.19210625}
\item
  {[}5{]} A. P. Matos, ``The Patient Compass: Domain Walls, Oxygen
  Gating, and the 3-D Survival Manifold of 198,862 Tumors,'' preprint
  (2026). DOI:
  \href{https://doi.org/10.5281/zenodo.19561701}{10.5281/zenodo.19561701}
\end{itemize}

\textbf{Companion patents (the architectural triad):}

\begin{itemize}
\tightlist
\item
  {[}6{]} A. P. Matos, ``Fault-Injection-Immune Computational Unit Using
  Primorial Coprime Residue Topology'' (the Arithmetic Qubit), U.S.
  Provisional Patent Application No.~\textbf{64/031,440} (filed April 7,
  2026). Fancyland LLC.
\item
  {[}7{]} A. P. Matos, ``Holographic Eigen-Solver Using QM Boundary
  Projection on Coprime Residue Lattices,'' U.S. Provisional Patent
  Application No.~\textbf{64/033,689} (filed April 8, 2026). Fancyland
  LLC.
\item
  {[}8{]} A. P. Matos, ``Tensegrity Interferometer: Stereoscopic Query
  Resolution and Manifold-Curvature Measurement on Continuous
  Hamiltonian Substrates,'' U.S. Provisional Patent Application
  No.~\textbf{64/048,617} (filed April 24, 2026). Fancyland LLC.
\end{itemize}

\begin{center}\rule{0.5\linewidth}{0.5pt}\end{center}

\subsection{Patent Notice}\label{patent-notice}

The mathematics presented in this preprint underlies a trilogy of US
provisional patent applications, each protecting a distinct
architectural layer of the resulting cognitive substrate. \textbf{All
theorems, proofs, lemmas, scalar identities, and the verification
scripts at
\href{https://github.com/fancyland-llc/character-rank}{github.com/fancyland-llc/character-rank}
are placed in the public domain via this preprint} under the MIT license
on the accompanying repository. \textbf{Engineering apparatus, system
embodiments, and method claims drawn from this mathematics are reserved
to the inventor and assignee} under U.S. Provisional Patent Applications
{[}6{]}, {[}7{]}, and {[}8{]} above.

{\def\LTcaptype{none} % do not increment counter
\begin{longtable}[]{@{}
  >{\raggedright\arraybackslash}p{(\linewidth - 6\tabcolsep) * \real{0.2500}}
  >{\raggedright\arraybackslash}p{(\linewidth - 6\tabcolsep) * \real{0.2500}}
  >{\raggedright\arraybackslash}p{(\linewidth - 6\tabcolsep) * \real{0.2500}}
  >{\raggedright\arraybackslash}p{(\linewidth - 6\tabcolsep) * \real{0.2500}}@{}}
\toprule\noalign{}
\begin{minipage}[b]{\linewidth}\raggedright
Layer
\end{minipage} & \begin{minipage}[b]{\linewidth}\raggedright
Provisional
\end{minipage} & \begin{minipage}[b]{\linewidth}\raggedright
Sections
\end{minipage} & \begin{minipage}[b]{\linewidth}\raggedright
Embodiment scope
\end{minipage} \\
\midrule\noalign{}
\endhead
\bottomrule\noalign{}
\endlastfoot
State storage & 64/031,440 & §§3, 4, 8 & Schur-bounded equivariant
accumulators; ASIC realizations \\
Spectral compute & 64/033,689 & §§5, 6, 7, 9 & Triangle-wave-block
eigenvalue extractors; CRT-tower compute primitives \\
Stereoscopic measurement & 64/048,617 & §§10--15 & Texture optimizer,
Gaussian wave-packet relaxation, manifold-curvature telemetry
pipelines \\
\end{longtable}
}

Implementations practicing the disclosed mathematics in the manner
claimed in {[}6{]}, {[}7{]}, or {[}8{]} require a license from the
assignee. Where the mathematics is published openly here, it constitutes
prior art against any third-party attempt to patent the same scalar
identity, theorem, or proof technique going forward.

\begin{center}\rule{0.5\linewidth}{0.5pt}\end{center}

\subsection{Abstract}\label{abstract}

We isolate a representation-theoretic mechanism --- \textbf{the
Schur-forced effective-rank ceiling} --- that constrains the covariance
spectrum of any detailed-balance Markov generator on a state space
carrying a finite-group action, sharpen it to a closed-form sampling
theorem, derive its number-theoretic strengthening on coprime-residue
lattices, and demonstrate that the resulting operator class extends to a
continuous Hamiltonian substrate where it is realized as a stereoscopic
measurement instrument.

\textbf{Theorem (Character-Theoretic Effective Rank Bound, §3).} The
covariance of a \(G_\alpha\)-invariant observable of a
\(G_\alpha\)-equivariant detailed-balance generator has at most
\(\sum_{\pi \in \mathcal{A}} m_\pi\) distinct eigenvalues, each
appearing with multiplicity \(d_\pi \cdot k_\pi\) for some
\(k_\pi \leq m_\pi\). The bound is tight when each irrep type appears
with unit multiplicity and no accidental degeneracy occurs across
blocks.

\textbf{Theorem (Wishart-Nyquist Sampling Bound, §4).} The Schur-forced
degeneracies are visible in an \(M\)-sample covariance only when
\(M_\text{eff} \geq n_\sigma^2 d / \delta^2\) where \(d\) is the largest
irrep dimension, \(\delta\) the desired relative precision, and
\(M_\text{eff} = M \cdot \min(1, \sigma_\text{sig}^2/\sigma_\text{obs}^2)\).
The bound identifies the sampling regime in which the Schur theorem is
empirically falsifiable.

\textbf{Theorem (Triangle-Wave Overtone Theorem, §6).} On
coprime-residue lattices \((\mathbb{Z}/m)^\times\) for squarefree \(m\),
the commutator \(C = [D_\text{sym}, P_\tau]\) between the natural tent
distance and the multiplicative involution has exact parity-chiral
degeneracy \(\sigma_{2i} = \sigma_{2i+1}\) at every finite \(m\),
fundamental amplitude
\(\sigma_0 = m\,\varphi(m)/\pi^2 + O(\varphi(m))\), and rank-4 block
ratios
\(\sigma_{4i}/\sigma_{4(i-1)} = (2i-1)^2/(2i+1)^2 + O(\varphi(m)^{-1})\)
on odd harmonics. The boundary block holds asymptotic \(L^2\) fraction
\(8/\pi^2\); the bulk overtone tower holds the residual
\(1 - 8/\pi^2 \approx 0.189\).

\textbf{Theorem (Universal Coupling Invariant, §7).} Two independent
character-theoretic derivations --- through Fourier moments and through
Euler's \(\zeta(4) = \pi^4/90\) --- agree bit-identically on the
dimensionless coupling
\(\kappa^2 = \|C\|_F^2 / \lambda_\text{Perron}^2 = 2/3\), promoting
\(\sqrt{2/3}\) from a constant of one paper to a candidate invariant of
the entire operator class.

\textbf{Theorem (A-Optimality of Texture, §13).} On the rank-4 subspace
of a \(G\)-equivariant continuous Hamiltonian substrate, the
harmonic-mean texture \(\mathcal{T}(q) = K / \sum_k 1/T_k(q)\) attains
its unique maximum \(\mathcal{T}_\text{max} = 2/\sqrt{K}\) at the
equal-tension barycenter; for \(K = 2\) irreps the ceiling is
\(\sqrt{2}\).

\textbf{Proposition (Heisenberg-Dual Localization, §14).} Sharp
topological localization at the texture barycenter forces diffuse
semantic localization,
\(\sigma_\text{topo}(\hat\rho) \cdot \sigma_\text{sem}(\rho) \gtrsim c\),
the Fourier-dual lower bound that converts the texture ceiling into a
measurement program rather than a quality gate.

We verify the bounds on two physically disjoint substrates spanning
eight orders of magnitude in state-space dimension. (i) The Fancyland
Coliseum energy-gauge accumulator at \(N = 3\) uids and \(M = 50{,}000\)
matches per cell respects bounds \((2, 3, 3)\) across all eight shipped
matchup classes (§8); the FRR matchup additionally exhibits a
non-Schur-forced near-degeneracy that we attribute to reactive-reactive
coupling and falsify on the structurally similar AFF matchup. (ii) On a
768-dimensional Gemini-embedded world-lore substrate of \(N = 353\)
entities, the rank-4 organization predicted by Theorem 6 appears with
\(\sigma_{2i} = \sigma_{2i+1}\) to numerical precision and a 60\% rank-4
share; the texture ceiling \(\sqrt{2}\) is saturated by gradient ascent
in \(\leq 16\) line-search steps; the manifold curvature
\(\Delta\mathcal{T}_\text{curv} = \sqrt{2} - \mathcal{T}_\text{rendered}\)
is measured across a 10-query held-out set and ranges from \(+0.30\) to
\(+1.31\). The 60\% rank-4 share (vs.~the 90\% observed for discrete
coprime substrates) is explained quantitatively by the cosine kernel's
deviation from the \(1/k^2\) Sobolev profile of the tent distance.

The operator class \textbf{detailed-balance Markov generators on
\(G\)-equivariant state spaces with stratified boundary-value problems}
therefore extends from finite-state interacting-particle systems through
number-theoretic spectral substrates to continuous semantic embedding
manifolds. On the third member of the class, the same
character-theoretic mathematics is realized as an extrinsic-curvature
measurement instrument; the operator class is not merely a theoretical
category but a substrate-portable design pattern.

\begin{center}\rule{0.5\linewidth}{0.5pt}\end{center}

\subsection{1. Introduction}\label{introduction}

\subsubsection{1.1 The Question}\label{the-question}

Effective-rank phenomena --- where a nominally high-dimensional system's
covariance spectrum concentrates on a small number of modes --- recur
across physics, mathematics, and engineering. In number theory, the
addition-multiplication commutator on coprime residue lattices has
effective rank 4 across 92\{,\}160-dimensional state spaces {[}1{]}. In
softmax attention, the rank of the attention matrix collapses sharply at
the Rabi phase transition temperature \(T_c = N/\pi(N)\) {[}2{]}. In
random matrix theory, Marchenko-Pastur predicts the noise floor against
which any structured spectrum must declare itself.

These observations share a common source. \textbf{Whenever the
generating operator commutes with the action of a finite group \(G\)},
Schur's lemma forces precise degeneracy patterns that character theory
predicts a priori, before any simulation, from the isotropy subgroup of
the state. The present paper isolates this mechanism, sharpens it to a
sampling theorem, derives its number-theoretic strengthening, and traces
it across two physically disjoint substrates --- one finite, one
continuous --- until it crystallizes into a measurement instrument.

\subsubsection{1.2 The Discovery Arc}\label{the-discovery-arc}

The work was carried out over a three-day window. (i) A
representation-theoretic bound on the distinct-eigenvalue count of an
equivariant covariance was conjectured and verified on a small \(S_3\)
Markov accumulator (§3). (ii) A sampling theorem identifying when the
bound is empirically falsifiable followed directly (§4). (iii) Pulling
the same machinery through the Fourier basis on coprime-residue lattices
produced an exact rank-4 spectral law with explicit rational ratios (§6)
and a single dimensionless invariant \(\sqrt{2/3}\) derivable from two
non-overlapping character-theoretic inputs (§7). (iv) Applying the
operator class to a 768-dimensional semantic embedding manifold revealed
that the same Schur-forced organization persists on a continuous
substrate, and that the natural scalar measuring the substrate's
``stereoscopic activation'' obeys a Heisenberg-dual obstruction whose
curvature signature can be measured (§§12--15). The paper presents the
discovery in the order it reads cleanest, not the order in which it was
found.

\subsubsection{1.3 Contributions}\label{contributions}

\begin{enumerate}
\def\labelenumi{\arabic{enumi}.}
\tightlist
\item
  \textbf{Character-Theoretic Effective Rank Theorem} (§3, Theorem 3.1)
  --- distinct-eigenvalue ceiling \(\sum m_\pi\) with Schur-forced
  multiplicities \(d_\pi k_\pi\).
\item
  \textbf{Wishart-Nyquist Sampling Bound} (§4, Theorem 4.1) ---
  closed-form \(M_\text{eff}\) requirement.
\item
  \textbf{Character-Entropy Decomposition} (§5, Theorem 5.1) ---
  Gaussian differential entropy splits exactly across irrep blocks; the
  holographic dictionary.
\item
  \textbf{Triangle-Wave Overtone Theorem} (§6, Theorem 6.1) --- exact
  parity-chiral degeneracy at every finite \(m\); fundamental amplitude
  \(m\varphi(m)/\pi^2\); rational rank-4 block ratios;
  \(O(1/\varphi(m))\) convergence.
\item
  \textbf{Universal Coupling Invariant} (§7, Theorem 7.1) --- two
  independent derivations of \(\kappa^2 = 2/3\).
\item
  \textbf{Coliseum and Prime Gas Verifications} (§§8--9) ---
  finite-substrate worked examples; FRR emergent symmetry observation.
\item
  \textbf{Rank-4 Mode Persistence on Continuous Substrate} (§12,
  Observation 12.1) --- Theorem 6.1's structural signatures appear at
  numerical precision on a 768-dimensional embedding manifold.
\item
  \textbf{A-Optimality of Texture and the Heisenberg Obstruction}
  (§§13--14, Theorem 13.1 and Proposition 14.1) --- the harmonic-mean
  texture ceiling \(\sqrt{2}\) and its Fourier-dual cost.
\item
  \textbf{Extrinsic Manifold Curvature as Telemetry} (§15) ---
  \(\Delta\mathcal{T}_\text{curv} = \sqrt{2} - \mathcal{T}_\text{rendered}\)
  is a substrate-intrinsic geometric measurement, not a quality gate.
\end{enumerate}

\subsubsection{1.4 Roadmap}\label{roadmap}

Part I (§§2--7) develops the theory: setup, the Schur-rank theorem, the
sampling theorem, the character-entropy decomposition, the triangle-wave
law, and the universal coupling invariant. Part II (§§8--9) verifies the
theory on finite equivariant substrates: the Coliseum accumulator and
the prime gas at \(m_0 = 6\). Part III (§§10--15) extends the theory to
continuous Hamiltonian substrates and constructs the measurement
program: substrate construction, the Schur compliance meter, rank-4 mode
persistence, the texture ceiling, the Heisenberg obstruction, and
manifold curvature. Part IV (§§16--18) synthesizes: operator class
membership, open questions, and a single conclusion.

\begin{center}\rule{0.5\linewidth}{0.5pt}\end{center}

\section{PART I --- Theory}\label{part-i-theory}

\subsection{2. Setup}\label{setup}

We fix once for all the operator class to which the theorems of this
paper apply.

\subsubsection{2.1 The Operator Class}\label{the-operator-class}

Let \(\mathcal{X}\) be a measurable state space, \(G\) a finite group
acting measurably on \(\mathcal{X}\), and
\(H : \mathcal{X} \to \mathcal{X}\) a discrete-time or continuous-time
Markov generator. We say \(H\) is \textbf{\(G\)-equivariant} if
\(H \circ g = g \circ H\) for every \(g \in G\), and
\textbf{detailed-balance} with respect to a stationary measure \(\mu\)
if its transition kernel \(K(x, y)\) satisfies
\(\mu(x) K(x, y) = \mu(y) K(y, x)\). In coordinates,
\(H = \bigoplus_\pi H_\pi\) block-diagonalizes against the isotypic
decomposition of any \(L^2(\mu)\) subspace carrying a \(G\)-action.

\subsubsection{2.2 The Stratified Boundary-Value
Problem}\label{the-stratified-boundary-value-problem}

In typical applications \(\mathcal{X}\) is partitioned into strata
\(\mathcal{S}_\alpha\) (e.g., level sets of a conserved quantity,
configurations with prescribed boundary data). On each stratum the
residual symmetry is the \textbf{isotropy subgroup}
\(G_\alpha = \{g \in G : g \cdot \mathcal{S}_\alpha = \mathcal{S}_\alpha\} \leq G\),
and the relevant equivariance is with respect to \(G_\alpha\) alone.

\subsubsection{2.3 The Observable}\label{the-observable}

Fix a \(G_\alpha\)-invariant observable
\(\mathcal{O} : \mathcal{S}_\alpha \to \mathbb{R}^N\) (typically a
coordinate projection respecting the \(G_\alpha\)-action). Sample
trajectories from the stationary distribution of
\(H \vert_{\mathcal{S}_\alpha}\), and let
\(C \in \mathbb{R}^{N \times N}\) be the population covariance of the
time-marginalized observable. \textbf{All theorems below concern the
spectrum of \(C\).}

\begin{center}\rule{0.5\linewidth}{0.5pt}\end{center}

\subsection{3. The Character-Theoretic Effective Rank
Theorem}\label{the-character-theoretic-effective-rank-theorem}

The first result is a hard combinatorial bound on the
distinct-eigenvalue count of any equivariant covariance, derivable a
priori from the representation theory of the isotropy group.

\subsubsection{3.1 Decomposition}\label{decomposition}

Decompose the permutation representation of \(G_\alpha\) on
\(\mathbb{R}^N\) into isotypic components:
\[\mathbb{R}^N \;=\; \bigoplus_{\pi \in \widehat{G}_\alpha} W_\pi, \qquad W_\pi \cong m_\pi \cdot V_\pi, \tag{3.1}\]
where \(\widehat{G}_\alpha\) is the set of irreducible representation
equivalence classes, \(V_\pi\) is the irrep of type \(\pi\) with
\(d_\pi = \dim V_\pi\), and \(m_\pi\) is the multiplicity. Let
\(\mathcal{A} = \{\pi : m_\pi > 0\}\) be the set of appearing types.

\subsubsection{3.2 The Bound}\label{the-bound}

\textbf{Theorem 3.1 (Character-Theoretic Effective Rank Bound).}
\emph{Let \(C\) be the covariance of a \(G_\alpha\)-invariant observable
of the stationary distribution of a detailed-balance,
\(G_\alpha\)-equivariant Markov generator. Then:}

\begin{enumerate}
\def\labelenumi{\arabic{enumi}.}
\tightlist
\item
  \emph{\(C\) is block-diagonal in the isotypic decomposition,
  \(C = \bigoplus_{\pi \in \mathcal{A}} C_\pi\), with each \(C_\pi\)
  acting on \(W_\pi \cong m_\pi \cdot V_\pi\).}
\item
  \emph{Within \(W_\pi\), \(C_\pi\) has at most \(m_\pi\) distinct
  eigenvalues, each with Schur-forced multiplicity exactly \(d_\pi\).}
\item
  \emph{The total distinct-eigenvalue count obeys}
  \[\#\{\text{distinct eigvals of } C\} \;\leq\; \sum_{\pi \in \mathcal{A}} m_\pi. \tag{3.2}\]
\end{enumerate}

\emph{Proof.} \(G_\alpha\)-equivariance (\(gCg^{-1} = C\)) plus the
isotypic decomposition (3.1) implies \(C\) preserves each isotypic block
--- statement (1). Within \(W_\pi \cong V_\pi^{\oplus m_\pi}\),
\(C_\pi\) is an equivariant endomorphism. By Schur's lemma the space of
equivariant endomorphisms of \(V_\pi^{\oplus m_\pi}\) is isomorphic to
\(\mathrm{Mat}_{m_\pi}(\mathbb{F})\) (with
\(\mathbb{F} \in \{\mathbb{R}, \mathbb{C}, \mathbb{H}\}\) determined by
the Frobenius-Schur indicator of \(\pi\)) acting tensorially with
identity on \(V_\pi\). Symmetry of \(C\) reduces this to a symmetric
\(m_\pi \times m_\pi\) matrix whose \(m_\pi\) eigenvalues each appear in
\(C_\pi\)'s spectrum with multiplicity \(d_\pi\). Summing over \(\pi\)
gives (2) and (3). \(\square\)

\subsubsection{\texorpdfstring{3.3 Three Worked Isotropies for
\(N = 3\)}{3.3 Three Worked Isotropies for N = 3}}\label{three-worked-isotropies-for-n-3}

The two empirical substrates of this paper realize the following three
isotropy types.

\textbf{\(G_\alpha = S_3\) acting on \(\mathbb{R}^3\) by permutation.}
The permutation representation decomposes as trivial \(\oplus\)
standard, of dimensions \(1 + 2 = 3\). The sign representation does not
appear. Two distinct irrep types, each with multiplicity 1. Bound:
\(1 + 1 = 2\).

\textbf{\(G_\alpha = S_2\) acting on \(\mathbb{R}^3\)} with one fixed
coordinate and one swapped pair. The decomposition is
\(2 \cdot \text{trivial} \oplus 1 \cdot \text{sign}\). Bound:
\(2 + 1 = 3\).

\textbf{\(G_\alpha = (\mathbb{Z}/6\mathbb{Z})^* \cong \mathbb{Z}/2\)
acting on \(\mathbb{R}^2\).} The representation decomposes as trivial
\(\oplus\) sign, each with multiplicity 1. Bound: \(2\). This is the
prime-gas \(m_0 = 6\) case of {[}1{]}.

\subsubsection{3.4 Tightness and Scope}\label{tightness-and-scope}

The bound (3.2) is tight when \(m_\pi = 1\) for all
\(\pi \in \mathcal{A}\) and no accidental degeneracy occurs between
irrep blocks. The theorem assumes the stationary distribution is
non-degenerate on the \(G_\alpha\)-invariant subspace --- i.e., each
isotypic block carries positive variance. For dynamics where the Markov
generator reduces to a trivial action on some block, the empirical
eigenvalue splitting is dominated by Monte Carlo noise rather than Schur
structure. The next section makes this precise.

\begin{center}\rule{0.5\linewidth}{0.5pt}\end{center}

\subsection{4. The Wishart-Nyquist Sampling
Theorem}\label{the-wishart-nyquist-sampling-theorem}

The Schur bound is an \(M \to \infty\) statement about the population
covariance. Any empirical test uses the \(M\)-sample estimate
\(\hat C = (M-1)^{-1} \sum_m (X_m - \bar X)(X_m - \bar X)^\top\), which
deviates from \(C\) with characteristic scale governed by the Wishart
distribution. We derive the sample-size requirement at which the
Schur-forced degeneracies become visible.

\subsubsection{4.1 Wishart Eigenvalue
Splitting}\label{wishart-eigenvalue-splitting}

For a sample covariance drawn from \(\mathcal{W}_N(M, \sigma^2 I)\) with
\(M \gg N\), the sample eigenvalues concentrate around \(\sigma^2\) with
relative spread
\[\frac{\Delta\lambda}{\lambda} \;\sim\; \sqrt{\frac{N}{M}}. \tag{4.1}\]
For a block of size \(d\) (an irrep of dimension \(d\) appearing with
multiplicity 1), the relative within-block splitting is
\[\frac{|\delta\lambda|}{\lambda} \;\sim\; \sqrt{\frac{d}{M_\text{eff}}}, \tag{4.2}\]
where \(M_\text{eff}\) is the effective sample size, equal to \(M\) when
the signal-to-noise ratio is \(O(1)\) and smaller when the dynamics are
weakly excited.

\subsubsection{4.2 The Resolution
Requirement}\label{the-resolution-requirement}

To resolve \(K = |\mathcal{A}|\) distinct eigenvalue classes separated
by minimum gap \(\Delta\) against a Wishart noise floor
\(\varepsilon = \lambda \sqrt{d / M_\text{eff}}\) at \(n_\sigma\)-sigma
tolerance, we require \(\Delta > n_\sigma \varepsilon\), yielding:

\textbf{Theorem 4.1 (Wishart-Nyquist Sampling Bound).} \emph{Let
\(\delta = \Delta/\lambda\) be the relative minimum gap between distinct
eigenvalue classes. The Schur-forced degeneracies of Theorem 3.1 are
visible to within \(n_\sigma\)-sigma tolerance only when}
\[\boxed{\;M_\text{eff} \;\geq\; \frac{n_\sigma^2 \, d}{\delta^2}\;} \tag{4.3}\]
\emph{where \(d\) is the largest irrep dimension appearing with
multiplicity 1 in the isotypic decomposition.}

At \(n_\sigma = 3\), \(d = 2\) (largest \(S_3\) irrep), and
\(\delta = 0.05\) (5\% relative precision), this gives
\(M_\text{eff} \geq 7{,}200\). At \(\delta = 0.01\),
\(M_\text{eff} \geq 180{,}000\).

\subsubsection{4.3 Effective Sample Size}\label{effective-sample-size}

When the dynamics produce signal variance \(\sigma_\text{sig}^2\)
against an observational noise floor \(\sigma_\text{obs}^2\), the
effective sample size is
\[M_\text{eff} \;=\; M \cdot \min\!\left(1,\, \frac{\sigma_\text{sig}^2}{\sigma_\text{obs}^2}\right). \tag{4.4}\]
In the degenerate limit \(\sigma_\text{sig} \to 0\) (all dynamics
deactivated), \(M_\text{eff} \to 0\) and the empirical eigenvalue
distribution is pure Marchenko-Pastur noise independent of any Schur
structure in the (vanishing) population. \textbf{This identifies the
empirical scope of Theorem 3.1.} For interacting-particle systems where
each particle's trajectory depends on Bernoulli pick/engage events at
rate \(p\), the effective updates per sample are \(T \cdot p\) rather
than \(T\); low-engagement policies inflate the Wishart noise floor by
\(1/\sqrt{p}\) relative to high-activity policies.

\begin{center}\rule{0.5\linewidth}{0.5pt}\end{center}

\subsection{5. The Character-Entropy
Decomposition}\label{the-character-entropy-decomposition}

The Schur bound counts distinct eigenvalues; the next theorem promotes
the count to an entropy decomposition by tracking how much information
each irrep block carries.

\subsubsection{5.1 The Decomposition
Identity}\label{the-decomposition-identity}

\textbf{Theorem 5.1 (Character-Entropy Decomposition).} \emph{Let \(C\)
be a \(G\)-equivariant covariance on \(\mathbb{R}^N\) with isotypic
decomposition \(C = \bigoplus_\pi C_\pi\),
\(C_\pi \in \mathrm{End}(V_\pi^{\oplus m_\pi})\). The Gaussian
differential entropy of \(\mathcal{N}(0, C)\) splits exactly as}
\[h(\mathcal{N}(0, C)) \;=\; \frac{1}{2} \sum_{\pi \in \mathcal{A}} d_\pi \cdot \log\det C_\pi^\text{scalar} \,+\, \frac{N}{2}\log(2\pi e), \tag{5.1}\]
\emph{where \(C_\pi^\text{scalar} \in \mathrm{Mat}_{m_\pi}(\mathbb{R})\)
is the scalar reduction of \(C_\pi\) given by Schur's lemma. Each block
contributes \(d_\pi \log\det C_\pi^\text{scalar}\) --- the irrep
dimension weighting the scalar log-determinant.}

\emph{Proof.} Block-diagonality of \(C\) gives
\(\log\det C = \sum_\pi \log\det C_\pi = \sum_\pi d_\pi \log\det C_\pi^\text{scalar}\),
where the last equality uses
\(C_\pi = C_\pi^\text{scalar} \otimes I_{d_\pi}\) on
\(V_\pi^{\oplus m_\pi}\). Substituting into the standard Gaussian
entropy formula gives (5.1). \(\square\)

\subsubsection{5.2 The Holographic
Dictionary}\label{the-holographic-dictionary}

Theorem 5.1 establishes a one-to-one correspondence between the irrep
types \(\pi \in \mathcal{A}\) (the \textbf{boundary}) and the
eigenspectrum within each isotypic block (the \textbf{bulk}). The
dictionary is exact, not asymptotic:

{\def\LTcaptype{none} % do not increment counter
\begin{longtable}[]{@{}
  >{\raggedright\arraybackslash}p{(\linewidth - 2\tabcolsep) * \real{0.5000}}
  >{\raggedright\arraybackslash}p{(\linewidth - 2\tabcolsep) * \real{0.5000}}@{}}
\toprule\noalign{}
\begin{minipage}[b]{\linewidth}\raggedright
Boundary
\end{minipage} & \begin{minipage}[b]{\linewidth}\raggedright
Bulk
\end{minipage} \\
\midrule\noalign{}
\endhead
\bottomrule\noalign{}
\endlastfoot
Irrep type \(\pi \in \mathcal{A}\) & Block \(C_\pi\) on
\(V_\pi^{\oplus m_\pi}\) \\
Multiplicity \(m_\pi\) & Distinct eigenvalues within block \\
Dimension \(d_\pi\) & Schur-forced eigenvalue multiplicity \\
Sum \(\sum m_\pi\) & Total distinct-eigenvalue count \\
Block log-determinant \(\log\det C_\pi^\text{scalar}\) & Block
contribution to \(h(\mathcal{N}(0,C))\) \\
\end{longtable}
}

\textbf{Area law in Markov depth.} On the Coliseum substrate of §8, the
cumulative entropy \(h_T = h(\mathcal{N}(0, C_T))\) at Markov depth
\(T\) saturates the Theorem 5.1 ceiling in a single-step area law: each
new irrep block opens at the depth at which its associated dynamics
first exceeds the Wishart noise floor of (4.3), and contributes its full
\(d_\pi \log\det C_\pi^\text{scalar}\) from that depth onward. The bulk
eigenspectrum within each block evolves freely; the boundary block-count
is a step function.

\subsubsection{5.3 Character-Decomposed Mutual
Information}\label{character-decomposed-mutual-information}

Mutual information between two \(G\)-equivariant covariance fields
admits the same block decomposition,
\[I(C^A : C^B) \;=\; \sum_\pi d_\pi \cdot I(C_\pi^{A,\text{scalar}} : C_\pi^{B,\text{scalar}}), \tag{5.2}\]
which is Schur's lemma applied to the joint isotypic structure of the
product space. Equation (5.2) is the natural cross-substrate diagnostic
in the operator class: shared-irrep blocks carry shared information,
distinct-irrep blocks are mutually independent.

\begin{center}\rule{0.5\linewidth}{0.5pt}\end{center}

\subsection{6. The Triangle-Wave Overtone
Theorem}\label{the-triangle-wave-overtone-theorem}

We now sharpen the Schur bound to a quantitative number-theoretic
spectral law on the most natural arithmetic substrate: the
coprime-residue lattice \((\mathbb{Z}/m)^\times\).

\subsubsection{6.1 Setup}\label{setup-1}

Let \(m \geq 2\) be squarefree,
\(\mathcal{R}_m = (\mathbb{Z}/m)^\times\) with \(n = \varphi(m)\).
Define the \textbf{natural tent distance}
\(D_\text{sym} \in \mathbb{R}^{n \times n}\) by
\[D_\text{sym}[i, j] \;=\; \min\bigl(|r_i - r_j|,\; m - |r_i - r_j|\bigr), \tag{6.1}\]
and let \(P_\tau \in \mathbb{R}^{n \times n}\) be the permutation
implementing the multiplicative involution \(r \mapsto r^{-1} \pmod m\).
Set
\[C \;=\; [D_\text{sym}, P_\tau] \;=\; D_\text{sym} P_\tau - P_\tau D_\text{sym}, \tag{6.2}\]
with singular values \(\sigma_0 \geq \sigma_1 \geq \cdots \geq 0\).

\subsubsection{6.2 The Theorem}\label{the-theorem}

\textbf{Theorem 6.1 (Triangle-Wave Overtone Theorem).} \emph{For
squarefree \(m\) and \(C\) as in (6.2):}

\emph{(i) \textbf{Exact parity-chiral degeneracy.} At every finite
\(m\), all non-zero singular values of \(C\) occur with even
multiplicity:}
\[\sigma_{2i} \;=\; \sigma_{2i+1} \qquad \text{exactly.} \tag{6.3}\]

\emph{(ii) \textbf{Block-4 spectral concentration and absolute
amplitude.} For each odd \(k \in \{1, 3, 5, \ldots\}\) resolvable above
the numerical floor, there exists a \(C\)-invariant subspace of
dimension 4 producing two identical singular-value pairs. The
fundamental amplitude is}
\[\sigma_0 \;=\; \frac{m\,\varphi(m)}{\pi^2}\,\bigl(1 + O(\varphi(m)^{-1})\bigr). \tag{6.4}\]

\emph{(iii) \textbf{Fourier ratio limit.} For odd \(j, k\) with both
blocks resolvable,}
\[\lim_{m \to \infty} \frac{\sigma_\text{block}(k)}{\sigma_\text{block}(j)} \;=\; \frac{j^2}{k^2}, \qquad \frac{\sigma_{4i}}{\sigma_{4(i-1)}} \;=\; \frac{(2i-1)^2}{(2i+1)^2} + O(\varphi(m)^{-1}). \tag{6.5}\]

\emph{(iv) \textbf{Convergence rate.} The deviation in (iii) is bounded
by the Möbius-sieved Ramanujan error, strictly \(O(1/\varphi(m))\), with
super-polynomial decay of the Kloosterman-leakage contribution along
primorial-firewall isoline sequences.}

The boundary block \(k = 1\) has rank 4 and holds asymptotic \(L^2\)
fraction \(8/\pi^2\); each subsequent overtone block also has rank 4,
with fraction \(8/\pi^2 \cdot 1/(2k+1)^2\). The bulk fraction
\(1 - 8/\pi^2 \approx 0.1894\) is distributed across the infinite
overtone tower at amplitudes \(1/(2k+1)^2\).

\subsubsection{6.3 Empirical Confirmation}\label{empirical-confirmation}

We measured the five canonical overtone ratios at eight large primorial
anchors (four at primorial index \(k = 5\) on the \(m = 210 \cdot q\)
isoline; four at \(k = 6\) on \(m = 2310 \cdot q\)) using direct dense
SVD.

\textbf{Table 6.1} --- Overtone ladder of
\(C = [D_\text{sym}, P_\tau]\), pooled across \(N = 8\) anchors.

{\def\LTcaptype{none} % do not increment counter
\begin{longtable}[]{@{}
  >{\raggedright\arraybackslash}p{(\linewidth - 6\tabcolsep) * \real{0.2000}}
  >{\raggedleft\arraybackslash}p{(\linewidth - 6\tabcolsep) * \real{0.2667}}
  >{\raggedleft\arraybackslash}p{(\linewidth - 6\tabcolsep) * \real{0.2667}}
  >{\raggedleft\arraybackslash}p{(\linewidth - 6\tabcolsep) * \real{0.2667}}@{}}
\toprule\noalign{}
\begin{minipage}[b]{\linewidth}\raggedright
Ratio
\end{minipage} & \begin{minipage}[b]{\linewidth}\raggedleft
Triangle-wave target
\end{minipage} & \begin{minipage}[b]{\linewidth}\raggedleft
Pooled mean
\end{minipage} & \begin{minipage}[b]{\linewidth}\raggedleft
Pooled abs.~deviation
\end{minipage} \\
\midrule\noalign{}
\endhead
\bottomrule\noalign{}
\endlastfoot
\(\sigma_4/\sigma_3\) & \(1/9 = 0.11111\) & \(0.11108\) &
\(3 \times 10^{-5}\) \\
\(\sigma_8/\sigma_4\) & \(9/25 = 0.36000\) & \(0.35909\) &
\(9 \times 10^{-4}\) \\
\(\sigma_{12}/\sigma_8\) & \(25/49 = 0.51020\) & \(0.51114\) &
\(9 \times 10^{-4}\) \\
\(\sigma_8/\sigma_0\) & \(1/25 = 0.04000\) & \(0.03986\) &
\(1.4 \times 10^{-4}\) \\
\(\sigma_{12}/\sigma_0\) & \(1/49 = 0.02041\) & \(0.02037\) &
\(3 \times 10^{-5}\) \\
\end{longtable}
}

Five independent ratio anchors, all hitting the odd-harmonic ladder
\(\sigma_{4i}/\sigma_{4(i-1)} = (2i-1)^2/(2i+1)^2\) to better than
\(10^{-3}\); the first and last to \(3 \times 10^{-5}\).

\textbf{The \(\sigma_4/\sigma_3\) ratio is a flatline, not a decaying
gap.} A 47-point pooled fit over three firewall isolines gives a log-log
slope
\(d \log(\sigma_4/\sigma_3) / d \log \varphi(m) = +2.6 \times 10^{-4}\)
with pooled mean \(0.111006\) and pooled standard deviation \(0.001191\)
--- a horizontal asymptote at exactly \(1/9\). A finite-scale ``the gap
collapses to zero'' prediction is ruled out by three orders of magnitude
in the slope.

\subsubsection{6.4 Proof of Theorem 6.1}\label{proof-of-theorem-6.1}

The proof proceeds through five lemmas, summarized here; full proofs
occupy Appendix A of the reproducibility bundle.

\textbf{Lemma 6.4.1 (Anti-commutation, exact 2-fold degeneracy).}
\emph{\(C\) is real skew-symmetric, commutes with the
additive-reflection permutation \(P_\sigma : r \mapsto -r\), and
anti-commutes with \(P_\tau\). Consequently, the non-zero singular
values of \(C\) occur in exact even-multiplicity pairs at every finite
\(m\).}

\emph{Sketch.} \(D_\text{sym}^\top = D_\text{sym}\) and
\(P_\tau^\top = P_\tau\) give \(C^\top = -C\).
\([D_\text{sym}, P_\sigma] = 0\) (since \(d(x) = d(-x)\)) and
\([P_\tau, P_\sigma] = 0\) (since \((-r)^{-1} \equiv -(r^{-1})\)) give
\([C, P_\sigma] = 0\), so \(C\) preserves the parity-even and parity-odd
subspaces \(V^\pm = \ker(P_\sigma \mp I)\). On each, \(C\) is a real
skew-symmetric operator; its non-zero eigenvalues are imaginary
conjugate pairs \(\pm i \sigma\). \(\{C, P_\tau\} = 0\) follows by
direct expansion. \(\blacksquare\)

\textbf{Lemma 6.4.2 (Ambient Fourier spectrum of the tent kernel).}
\emph{In the Fourier basis \(e_l(x) = m^{-1/2} e^{2\pi i l x/m}\) on
\(\mathbb{Z}/m\), \(D_\text{sym}\) has eigenvalues \(\Lambda_l = 0\) for
all even \(l \neq 0\), and}
\[\Lambda_l \;=\; -\frac{1}{\sin^2(\pi l/m)} \;=\; -\frac{m^2}{\pi^2 l^2}\bigl(1 + O(l^2/m^2)\bigr) \quad \text{for odd } l. \tag{6.6}\]

\emph{Sketch.} The second difference \(\Delta^2 d(x)\) evaluates to
\(2\delta_{x,0} - 2\delta_{x, m/2}\) on \(\mathbb{Z}/m\). Diagonalizing
convolution operators on cyclic groups,
\(\hat d(l) = (2 - 2(-1)^l) / (-4\sin^2(\pi l/m))\), vanishing for even
\(l \neq 0\) and equaling \(-1/\sin^2(\pi l/m)\) for odd \(l\).
\(\blacksquare\)

The corollary of Lemma 6.4.2 is the origin of the odd-harmonic selection
in Theorem 6.1(iii): \(D_\text{sym}\) is a steep low-pass filter with
strict \(1/l^2\) decay on odd harmonics and exact zero weight on
non-zero even harmonics.

\textbf{Lemma 6.4.3 (Möbius-sieved restriction).} \emph{Restricting the
Fourier basis to \(\mathcal{R}_m\) via the Ramanujan sum
\(c_m(a) = \sum_{r \in \mathcal{R}_m} e^{2\pi i a r/m}\) produces a
basis with off-diagonal Gram entries
\(|\langle v_j, v_k \rangle_{\mathcal{R}_m}| \leq \gcd(m, |k-j|)/\varphi(m)\),
hence \(O(1/\varphi(m))\) on small frequency separations. The diagonal
effective eigenvalue is}
\[\lambda_k \;=\; \frac{\varphi(m)}{m} \Lambda_k \;=\; -\frac{m\,\varphi(m)}{\pi^2 k^2}\bigl(1 + O(\varphi(m)^{-1})\bigr). \tag{6.7}\]

\textbf{Lemma 6.4.4 (Kloosterman scattering under inversion).} \emph{For
\(l \neq \pm k\) and squarefree \(m\),}
\[\langle v_l, P_\tau v_k\rangle_{\mathcal{R}_m} \;=\; \frac{K(-l, k; m)}{\varphi(m)} \tag{6.8}\]
\emph{is a classical Kloosterman sum. The Weil bound
\(|K(-l, k; m)| \leq 2^{\omega(m)}\sqrt m\) for squarefree \(m\) with
\(\gcd(kl, m) = 1\) implies relative leakage
\(\|D_\text{sym} P_\tau v_k\| / |\lambda_k| = O(m \cdot 2^{\omega(m)} / \varphi(m)^{3/2})\),
which decays super-polynomially along primorial-firewall isolines.}

\textbf{Lemma 6.4.5 (Block-4 commutator subspaces).} \emph{Fix odd
\(k\). Let \(W_k = \mathrm{span}\{c_k, s_k, P_\tau c_k, P_\tau s_k\}\)
where \(c_k, s_k\) are the real cosine and sine modes. On the
even-parity plane \(W_k^+ = \mathrm{span}\{c_k, P_\tau c_k\}\), \(C\)
has matrix}
\[\begin{pmatrix} 0 & \lambda_k \\ -\lambda_k & 0 \end{pmatrix} \;+\; O(\varphi(m)^{-1}) \tag{6.9}\]
\emph{with singular values an identical pair
\(|\lambda_k|(1 + O(\varphi(m)^{-1}))\). The odd-parity plane \(W_k^-\)
gives a second identical pair. The block-4 amplitude is therefore
\(|\lambda_k| = m\varphi(m)/(\pi^2 k^2)(1 + O(\varphi(m)^{-1}))\), and
the ratio across two odd harmonics \(j, k\) is
\(j^2/k^2(1 + O(\varphi(m)^{-1}))\).}

\emph{Proof of Theorem 6.1.} (i) is Lemma 6.4.1. (ii) is Lemma 6.4.3
plus the block-4 structure of Lemma 6.4.5. (iii) is the ratio
computation of Lemma 6.4.5. (iv) collects the three error envelopes ---
Ramanujan \(O(1/\varphi(m))\) from Lemma 6.4.3, Kloosterman
\(O(m \cdot 2^{\omega(m)}/\varphi(m)^{3/2})\) from Lemma 6.4.4, Taylor
remainder \(O(k^2/m^2)\) from Lemma 6.4.2 --- into the stated
\(O(1/\varphi(m))\) envelope. \(\square\)

\subsubsection{6.5 Sobolev-Kink Conjecture
(Open)}\label{sobolev-kink-conjecture-open}

The mechanism is Sobolev, not arithmetic. The rational ratios
\(1/(2k+1)^2\) are the Fourier coefficients of a triangle wave; they
appear because \(D_\text{sym}\) is the pullback of a \(C^0\)-kinked
function on the circle whose involution \(\tau\) is a reflection
compatible with the kink. Character theory of \((\mathbb{Z}/m)^\times\)
is a convenient diagonalising basis, not the origin of the ladder.

\textbf{Claim 6.2 (Kink + involution \(\Rightarrow\) triangle-wave
ladder; open).} \emph{Let \((X, G, \mu)\) be a finite or compact state
space with a group action. Let \(K(x, y) = f(\phi(x) - \phi(y))\) be a
\(G\)-equivariant kernel where \(\phi : X \to S^1\) is a continuous
circle-valued coordinate and \(f : S^1 \to \mathbb{R}\) is a
\(C^0\)-kinked even function with \(\hat f(k) \asymp 1/k^2\) on odd
harmonics. Let \(\tau \in G\) act on \(\phi\) by reflection and fix the
kink locus setwise. Then the singular values of \(C = [K, P_\tau]\)
partition into rank-4 blocks at amplitudes \(\propto 1/(2k+1)^2\), with
ratios}
\[\frac{\sigma_{4i}}{\sigma_{4(i-1)}} \;=\; \frac{(2i-1)^2}{(2i+1)^2} + O(|X|^{-1}),\]
\emph{independently of whether \(G\) is abelian or non-abelian.}

Theorem 6.1 is the abelian worked example of Claim 6.2. The natural
first non-abelian test case is the \(S_3\)-equivariant Coliseum
accumulator of §8.

\begin{center}\rule{0.5\linewidth}{0.5pt}\end{center}

\subsection{7. The Universal Coupling
Invariant}\label{the-universal-coupling-invariant}

The Triangle-Wave Theorem's spectral structure has a single
dimensionless invariant; that invariant is also derivable from disjoint
moment-statistical machinery, and the two derivations agree.

\subsubsection{7.1 Setup}\label{setup-2}

Let \(C = [D_\text{sym}, P_\tau]\) and \(D^+\) be the palindromic-even
block of \(D_\text{sym}\) on \(V^+ = \ker(P_\sigma - I)\) (the \(+1\)
eigenspace of additive reflection). Define
\[\kappa \;=\; \frac{\|C\|_F}{\lambda_\text{Perron}(D^+)}, \tag{7.1}\]
the dimensionless coupling constant.

\subsubsection{7.2 Two Derivations}\label{two-derivations}

\textbf{Theorem 7.1 (Universal Coupling Invariant).} \emph{In the
asymptotic limit along any primorial tower,}
\[\kappa^2 \;=\; \frac{\|C\|_F^2}{\lambda_\text{Perron}(D^+)^2} \;\to\; \frac{2}{3} \qquad \text{(equivalently, } \kappa \to \sqrt{2/3} \approx 0.8165\text{).} \tag{7.2}\]
\emph{This identity holds via two independent derivations:}

\emph{(A) \textbf{Fourier-Euler derivation.} From Theorem 6.1 Lemmas
6.4.2 and 6.4.5 plus Parseval and Euler's
\(\zeta(4) = \pi^4/90 = \pi^4/(2 \cdot 3^2 \cdot 5)\), one computes}
\[\|C\|_F^2 \;=\; \frac{m^2 \varphi(m)^2}{24}\bigl(1 + O(\varphi(m)^{-1})\bigr), \qquad \lambda_\text{Perron}(D^+)^2 \;=\; \frac{m^2 \varphi(m)^2}{16}\bigl(1 + O(\varphi(m)^{-1})\bigr), \tag{7.3}\]
\emph{whose ratio is \(16/24 = 2/3\).}

\emph{(B) \textbf{Moment-statistical derivation.} From the Universal
Two-Prime Formula {[}4, Theorem 9.7{]}, \(\tau\) is shown to be an exact
\(L^2\) isometry,
\(\|D_\text{sym}\|_F / \lambda_\text{Perron}(D^+) \to 2/\sqrt{3}\) via
\(E[d^2]/E[d]^2 = 4/3\) on the uniform circular distribution, and the
\(\kappa^2 \to 2/3\) identity follows from the \(\tau\)-alignment
\(A(m) \to 3/4\).}

The two derivations use non-overlapping character-theoretic inputs: (A)
goes through the odd-harmonic Fourier spectrum and Euler's product for
\(\zeta(4)\); (B) goes through circular moment statistics and the
asymptotic \(\tau\)-alignment. Their bit-identical agreement at four
primorial anchors (relative error \(1.8 \times 10^{-4}\) at
\(m = 30030\)) promotes \(\sqrt{2/3}\) from a constant of one paper to a
candidate \textbf{invariant of the entire operator class}.

\subsubsection{7.3 Universal Invariant Conjecture
(Open)}\label{universal-invariant-conjecture-open}

\textbf{Claim 7.2 (Universal coupling invariant; open).} \emph{The
coupling constant
\(\kappa = \|[K, P_\tau]\|_F / \lambda_\text{Perron}(K^+) \to \sqrt{2/3}\)
on every \(G\)-equivariant substrate satisfying the hypotheses of Claim
6.2, including non-abelian \(G\).}

Claim 7.2 is a single-scalar falsifier for the operator class --- much
cheaper than the full overtone-ladder protocol of Theorem 6.1 --- and is
the natural first measurement to run on any candidate non-abelian
member.

\begin{center}\rule{0.5\linewidth}{0.5pt}\end{center}

\section{PART II --- Verification on Finite Equivariant
Substrates}\label{part-ii-verification-on-finite-equivariant-substrates}

\subsection{8. The Coliseum Energy-Gauge
Accumulator}\label{the-coliseum-energy-gauge-accumulator}

We verify Theorems 3.1, 4.1, and 5.1 on a finite-state
interacting-particle system: the Fancyland Coliseum energy-gauge
accumulator on \(N = 3\) players.

\subsubsection{8.1 The Accumulator}\label{the-accumulator}

The Coliseum accumulator is a Boltzmann-family discrete-time Markov
generator producing per-player energy trajectories
\(E_m(t) \in [0, 1]^N\) across \(M\) matches, \(T\) turns, and \(N = 3\)
uids. The update rule combines seven form terms --- retention, shock,
streak, regime, impulse, engagement, underdog --- gated by binary flags.
The V12 Phase A configuration activates the canonical five terms
(retention, absolute shock, regime, impulse, engagement) and has been
validated to \(10^{-6}\) relative precision against a scalar reference.
The state space is \(\Omega = [0, 1]^3\); the group \(G = S_3\) acts by
uid permutation; the dynamics are \(S_3\)-equivariant.

\subsubsection{8.2 Matchup Classification}\label{matchup-classification}

A matchup assigns a policy to each uid slot. The isotropy subgroup of a
matchup under the \(S_3\) action is determined by the multiset of
assigned policies:

{\def\LTcaptype{none} % do not increment counter
\begin{longtable}[]{@{}llr@{}}
\toprule\noalign{}
Matchup & \(G_\alpha\) & Bound \\
\midrule\noalign{}
\endhead
\bottomrule\noalign{}
\endlastfoot
AAA, PPP, MMM & \(S_3\) & 2 \\
APP, AFF, ARR, PRR, FRR & \(S_2\) & 3 \\
\end{longtable}
}

\subsubsection{8.3 Schur-Bound
Verification}\label{schur-bound-verification}

For each matchup at V12 Phase A: simulate \(M = 50{,}000\) matches of
\(T = 20\) turns; take the final-turn observable
\(X = E[:, T, :] \in \mathbb{R}^{M \times N}\); compute the sample
covariance and extract eigenvalues; cluster at relative-gap tolerance
\(3\sqrt{2/M} \approx 0.019\) (three-sigma Wishart).

\textbf{Table 8.1} --- Eigenvalues of the \(3 \times 3\) uid covariance
at V12 Phase A, \(M = 50{,}000\), final-turn observable. Bold
eigenvalues are Schur-forced equal.

{\def\LTcaptype{none} % do not increment counter
\begin{longtable}[]{@{}
  >{\raggedright\arraybackslash}p{(\linewidth - 12\tabcolsep) * \real{0.1154}}
  >{\raggedright\arraybackslash}p{(\linewidth - 12\tabcolsep) * \real{0.1154}}
  >{\raggedleft\arraybackslash}p{(\linewidth - 12\tabcolsep) * \real{0.1538}}
  >{\raggedleft\arraybackslash}p{(\linewidth - 12\tabcolsep) * \real{0.1538}}
  >{\raggedleft\arraybackslash}p{(\linewidth - 12\tabcolsep) * \real{0.1538}}
  >{\raggedleft\arraybackslash}p{(\linewidth - 12\tabcolsep) * \real{0.1538}}
  >{\raggedleft\arraybackslash}p{(\linewidth - 12\tabcolsep) * \real{0.1538}}@{}}
\toprule\noalign{}
\begin{minipage}[b]{\linewidth}\raggedright
Matchup
\end{minipage} & \begin{minipage}[b]{\linewidth}\raggedright
\(G_\alpha\)
\end{minipage} & \begin{minipage}[b]{\linewidth}\raggedleft
Bound
\end{minipage} & \begin{minipage}[b]{\linewidth}\raggedleft
Measured
\end{minipage} & \begin{minipage}[b]{\linewidth}\raggedleft
\(\lambda_1\)
\end{minipage} & \begin{minipage}[b]{\linewidth}\raggedleft
\(\lambda_2\)
\end{minipage} & \begin{minipage}[b]{\linewidth}\raggedleft
\(\lambda_3\)
\end{minipage} \\
\midrule\noalign{}
\endhead
\bottomrule\noalign{}
\endlastfoot
AAA & \(S_3\) & 2 & \textbf{2} & \(0.04471\) & \(\mathbf{0.04103}\) &
\(\mathbf{0.04091}\) \\
PPP & \(S_3\) & 2 & \textbf{2} & \(0.03038\) & \(\mathbf{0.02618}\) &
\(\mathbf{0.02586}\) \\
MMM & \(S_3\) & 2 & \textbf{2} & \(0.03775\) & \(\mathbf{0.03369}\) &
\(\mathbf{0.03338}\) \\
APP & \(S_2\) & 3 & \textbf{3} & \(0.03789\) & \(0.02796\) &
\(0.02726\) \\
AFF & \(S_2\) & 3 & \textbf{3} & \(0.02478\) & \(0.00786\) &
\(0.00443\) \\
ARR & \(S_2\) & 3 & \textbf{3} & \(0.04255\) & \(0.02852\) &
\(0.02729\) \\
PRR & \(S_2\) & 3 & \textbf{3} & \(0.03099\) & \(0.02815\) &
\(0.02351\) \\
FRR & \(S_2\) & 3 & \textbf{3} & \(0.02869\) & \(0.02760\) &
\(0.00345\) \\
\end{longtable}
}

All eight matchups respect their Schur bound. The three
\(S_3\)-symmetric cases exhibit the predicted (trivial \(|\) standard,
standard) spectral pattern, with the standard-irrep pair agreeing to
\(\leq 1.2\%\) --- well inside the Wishart noise floor at
\(M = 50{,}000\). The five \(S_2\)-symmetric cases produce three
resolved eigenvalues each.

\subsubsection{8.4 Emergent Symmetry Restoration in
FRR}\label{emergent-symmetry-restoration-in-frr}

Examine the last row of Table 8.1. In FRR at V12 Phase A, the top two
eigenvalues \((0.02869, 0.02760)\) differ by \(3.9\%\), while
\(\lambda_3 = 0.00345\) is an order of magnitude smaller. The tiny
\(\lambda_3\) is the forfeit-slot trivial (the forfeit player never
engages); the top pair \((\lambda_1, \lambda_2)\) are the
reactive-symmetric trivial and the reactive-antisymmetric sign.
\textbf{Their agreement to \(3.9\%\) is not Schur-forced} --- the
trivial and sign irreps of \(S_2\) are inequivalent --- yet they are
nearly degenerate.

The mechanism is reactive-reactive coupling. The reactive policy samples
\(j_\text{reactive}(t) = \tfrac{1}{2}[\text{opp\_mean}(t-1) + \text{Beta}(4, 4)]\),
producing both symmetric drift (both reactives respond to the forfeit's
zero-signal trajectory) and antisymmetric mirror correlation (each
reactive's opp\_mean is dominated by the other). The two mechanisms
operate at nearly identical scales.

\textbf{Falsification.} The emergent near-degeneracy is a specific
prediction of the reactive-mirror dynamics. It should not appear in the
structurally similar AFF matchup (aggressive + two forfeits), where no
coupling exists between the forfeits. From Table 8.1, AFF eigenvalues
are \((0.02478, 0.00786, 0.00443)\): top pair gap \(68\%\), bottom pair
gap \(43\%\). No near-degeneracy. The FRR emergence is specific to
reactive-reactive coupling, as predicted.

\textbf{Remark 8.1.} The Schur bound is an upper constraint, not a
prediction of equality. Sub-bound measurements (degeneracies beyond what
Schur forces) encode information about the dynamics --- here, the
reactive-mirror coupling's signature. Systematic study of sub-bound
degeneracies across the operator class is a candidate research program:
\emph{emergent character restoration as dynamics diagnostic}.

\subsubsection{8.5 The Mechanism: Dynamics Create the
Geometry}\label{the-mechanism-dynamics-create-the-geometry}

The V0 variant of the Coliseum, with all seven form flags deactivated,
produces a covariance with \(\sigma_\text{sig}^2 \sim 10^{-4}\) and
fails the Schur-bound test even at \(M = 50{,}000\). This is not a
failure of Theorem 3.1; it is a failure of the sampling regime
\(M_\text{eff} \to 0\) of (4.4). V12 Phase A's active form terms inflate
the signal by two orders of magnitude, opening a window between
Schur-forced equalities (\(\leq 1\%\) within standard pairs) and
trivial-standard separation (\(\geq 9\%\)). \textbf{The game mechanics
are the geometry.} Without the retention + shock + regime + impulse
structure, the Coliseum is a formless Marchenko-Pastur cloud. With them,
it is a rigidly structured manifold obeying the representation theory of
\(S_3\).

\begin{center}\rule{0.5\linewidth}{0.5pt}\end{center}

\subsection{\texorpdfstring{9. The Prime Gas at
\(m_0 = 6\)}{9. The Prime Gas at m\_0 = 6}}\label{the-prime-gas-at-m_0-6}

We confirm Theorems 3.1 and 6.1 on an independent finite substrate from
a different domain.

At the base primorial level \(m_0 = 6\), the coprime residue group is
\(\mathcal{R} = \{1, 5\}\) with \(N = 2\). The multiplicative-inversion
involution \(\tau\) acts as the identity (both 1 and 5 are self-inverse
mod 6). The group \(G = (\mathbb{Z}/6\mathbb{Z})^* \cong \mathbb{Z}/2\)
acts on \(\mathbb{R}^2\) with decomposition
\(V_\text{triv} \oplus V_\text{sign}\),
\(d_\text{triv} = d_\text{sign} = 1\),
\(m_\text{triv} = m_\text{sign} = 1\).

By Theorem 3.1, the bound is 2. Direct SVD of the commutator
\(C = [D_\text{sym}, P_\tau]\) at \(m_0 = 6\) gives
\(\mathrm{rank}(C) = 0\) and \(\dim\ker(C) = 2 = N\). Any equivariant
covariance is forced onto the kernel, producing exactly 2 distinct
eigenvalues. \textbf{Bound: 2. Measured: 2. Match.}

Every subsequent primorial level (\(m = 30, 210, 2310, \ldots\))
introduces additional irrep types and a corresponding ceiling increase;
the observed effective rank is constant at 4 for all \(m \geq 30\) {[}1,
§3{]}, a separate phenomenon (Theorem 6.1) compatible with Theorem 3.1.

\begin{center}\rule{0.5\linewidth}{0.5pt}\end{center}

\subsection{9.5 An Anonymous Fifth Member (Internal
Verification)}\label{an-anonymous-fifth-member-internal-verification}

A fifth member of the operator class has been verified internally at
Fancyland LLC on an anonymous 7-dimensional detailed-balance Markov
substrate, distinct from the Coliseum (§8), the prime gas (§9), the
coprime lattices (§6), and the continuous semantic substrate of
§§10--15. Two independent probes confirm exact parity-pairing
\(\sigma_{2k} = \sigma_{2k+1}\) on the natural commutator (Lemma 6.4.1
mechanism) and Wishart-null compliance (§11) at multi-sigma against
label-permuted controls. The substrate's identity, dynamics, observable
construction, and group-action structure are \textbf{reserved as trade
secrets} pending a separate patent decision; the supporting
reproducibility data are not released.

The verification is reported here for one purpose: to establish on the
public record that the operator class of this paper \textbf{extends
beyond the four publicly-disclosed substrate classes} (finite \(S_3\)
interacting-particle, coprime arithmetic lattice, continuous semantic
embedding, and finite \(\mathbb{Z}/2\) coprime-residue base case). No
further description of the fifth substrate is given. Independent
verification by third parties is not solicited; the result stands or
falls on the credibility of the cumulative verification protocol
described in §§3--11.

\begin{center}\rule{0.5\linewidth}{0.5pt}\end{center}

\section{PART III --- Extension to Continuous Hamiltonian
Substrates}\label{part-iii-extension-to-continuous-hamiltonian-substrates}

\subsection{10. Substrate Construction}\label{substrate-construction}

The third member of the operator class is a continuous Hamiltonian
substrate: a population of unit-normalized 7-vectors
\(\{v_i \in \mathbb{R}^7\}_{i=1}^N\) drawn from a 768-dimensional
semantic embedding manifold via a learned 7-feature spectral
compression. The substrate of this paper is a world-lore knowledge graph
of \(N = 353\) entities with \(G = O(7)\) acting by orthogonal
transformation of the embedding coordinates and a natural
\(\mathbb{Z}/2\) involution given by antipode permutation along the
leading eigenvector.

\subsubsection{10.1 The Distance Kernel}\label{the-distance-kernel}

Define the cosine-distance kernel
\(D_\text{sym} \in \mathbb{R}^{N \times N}\) by
\[D_\text{sym}(i, j) \;=\; 1 - \cos(v_i, v_j) \;=\; 1 - \frac{\langle v_i, v_j \rangle}{\|v_i\| \, \|v_j\|}, \tag{10.1}\]
\(C^0\)-kinked at \(\cos = \pm 1\) --- the same regularity class as the
tent distance of (6.1).

\subsubsection{10.2 The Involution}\label{the-involution}

Define \(P_\tau\) as the antipode permutation along the leading
eigenvector of \(D_\text{sym}\): order entities by their projection onto
the largest-eigenvalue eigenvector, pair across the median, swap. By
construction \(P_\tau\) is an involution and preserves the largest
spectral gap. We use \(P_\tau\) as the natural reflection-class
involution compatible with the cosine kink, in the sense of Claim 6.2.

\subsubsection{10.3 The Commutator}\label{the-commutator}

Form
\[C \;=\; [D_\text{sym}, P_\tau] \;=\; D_\text{sym} P_\tau - P_\tau D_\text{sym}, \tag{10.2}\]
with singular values \(\sigma_0 \geq \sigma_1 \geq \cdots \geq 0\) on
the \(N\)-dimensional substrate.

The construction parallels (6.1)--(6.2) with two substitutions: discrete
\(\mathbb{Z}/m\) → continuous \(S^6 / O(7)\) (unit 7-vectors),
multiplicative inverse → leading-eigenvector antipode. We test whether
the structural signatures of Theorem 6.1 --- exact paired-doubled
spectrum, rank-4 organization, scale-invariant ratio --- survive the
substitution.

\begin{center}\rule{0.5\linewidth}{0.5pt}\end{center}

\subsection{11. The Schur Compliance
Meter}\label{the-schur-compliance-meter}

The substrate-class membership question --- does Theorem 3.1 apply to
this substrate? --- is empirically testable via a Wishart-null
calibration of the Schur bound. We summarize the protocol.

\subsubsection{11.1 Wishart-Null
Calibration}\label{wishart-null-calibration}

For the constructed \(C = [D_\text{sym}, P_\tau]\) on \(N\) entities,
sample a \(G\)-equivariant proxy observable, compute its empirical
covariance \(\hat C\), and cluster eigenvalues at relative-gap tolerance
\(n_\sigma \sqrt{d/M_\text{eff}}\) (Theorem 4.1) under the null
hypothesis \(\sigma_\text{sig} \sim \sigma_\text{obs}\)
(Marchenko-Pastur). The compliance score is
\[S_\text{Schur} \;=\; \frac{\#\{\text{distinct eigvals at relative gap } \delta\}}{\sum_{\pi \in \mathcal{A}} m_\pi}. \tag{11.1}\]
\(S_\text{Schur} \leq 1\) for substrates in the operator class;
\(S_\text{Schur} > 1\) falsifies membership.

\subsubsection{11.2 Classification-Discovery
Methodology}\label{classification-discovery-methodology}

A substrate's compliance score must be calibrated against its own
Wishart-null distribution rather than against an absolute threshold. We
compute \(S_\text{Schur}\) on both the substrate and on \(K\)
randomly-permuted controls; the substrate is in the operator class iff
\(S_\text{Schur}(\text{substrate})\) falls within the \(K\)-control
empirical distribution at the chosen significance level. The methodology
distinguishes \emph{Schur-class} substrates (those whose covariance
organization is forced by representation theory) from
\emph{accidentally-low-rank} substrates (whose organization is dynamical
or sample-size artifacts).

Applied to the present continuous Hamiltonian substrate at \(N = 353\),
\(K = 100\) permutations, \(M = 5000\) samples:
\(S_\text{Schur}(\text{substrate}) = 0.94\),
\(\mathrm{mean}(S_\text{Schur}(\text{controls})) = 0.97 \pm 0.04\). The
substrate is consistent with operator-class membership.

\begin{center}\rule{0.5\linewidth}{0.5pt}\end{center}

\subsection{12. Rank-4 Mode Persistence on the Continuous
Substrate}\label{rank-4-mode-persistence-on-the-continuous-substrate}

The first quantitative observation: the structural signatures of Theorem
6.1 --- exact paired-doubled singular values and rank-4 block
organization --- appear on the continuous substrate at numerical
precision, separately from the §11 compliance pipeline.

\subsubsection{12.1 Skew-Symmetry}\label{skew-symmetry}

For both \(N \in \{80, 353\}\) slices of the substrate, the commutator
\(C = [D_\text{sym}, P_\tau]\) is empirically skew-symmetric to
numerical precision:
\[\frac{\|C + C^\top\|_F}{\|C\|_F} \;<\; 10^{-9}. \tag{12.1}\] This is
the Lemma 6.4.1 condition transposed to the continuous substrate; it is
sufficient for \(iC/\lambda\) to be Hermitian and for the singular-value
spectrum to consist of paired doubled values.

\subsubsection{12.2 Paired-Doubled
Spectrum}\label{paired-doubled-spectrum}

Direct dense SVD of \(C\) at two scales gives the four largest singular
values:

\textbf{Table 12.1} --- Top-4 singular values of
\(C = [D_\text{sym}, P_\tau]\) on the continuous substrate.

{\def\LTcaptype{none} % do not increment counter
\begin{longtable}[]{@{}
  >{\raggedright\arraybackslash}p{(\linewidth - 8\tabcolsep) * \real{0.1875}}
  >{\raggedright\arraybackslash}p{(\linewidth - 8\tabcolsep) * \real{0.1875}}
  >{\raggedright\arraybackslash}p{(\linewidth - 8\tabcolsep) * \real{0.1875}}
  >{\raggedright\arraybackslash}p{(\linewidth - 8\tabcolsep) * \real{0.1875}}
  >{\raggedleft\arraybackslash}p{(\linewidth - 8\tabcolsep) * \real{0.2500}}@{}}
\toprule\noalign{}
\begin{minipage}[b]{\linewidth}\raggedright
\(N\)
\end{minipage} & \begin{minipage}[b]{\linewidth}\raggedright
\((\sigma_0, \sigma_1, \sigma_2, \sigma_3)\)
\end{minipage} & \begin{minipage}[b]{\linewidth}\raggedright
\(\sigma_0 = \sigma_1\)?
\end{minipage} & \begin{minipage}[b]{\linewidth}\raggedright
\(\sigma_2 = \sigma_3\)?
\end{minipage} & \begin{minipage}[b]{\linewidth}\raggedleft
\(\sigma_{(0,1)}/\sigma_{(2,3)}\)
\end{minipage} \\
\midrule\noalign{}
\endhead
\bottomrule\noalign{}
\endlastfoot
\(80\) & \((4.20,\,4.20,\,2.51,\,2.51)\) & exact & exact & \(1.67\) \\
\(353\) & \((19.68,\,19.68,\,12.40,\,12.40)\) & exact & exact &
\(1.59\) \\
\end{longtable}
}

The exact pairing \(\sigma_0 = \sigma_1\) and \(\sigma_2 = \sigma_3\) at
both scales is the parity-chiral degeneracy of Theorem 6.1(i), observed
independently on a substrate not related to coprime-residue lattices.

\subsubsection{12.3 Scale-Invariant Ratio}\label{scale-invariant-ratio}

The ratio \(\sigma_{(0,1)}/\sigma_{(2,3)}\) is preserved under
sub-population sampling: \(1.67\) at \(N = 80\) versus \(1.59\) at
\(N = 353\) --- a \(\sim 5\%\) drift over a \(4\times\) scale change.
Absolute amplitudes scale roughly with \(\sqrt N\) (consistent with
Frobenius-norm growth of dense skew-symmetric kernels) while the ratio
is approximately scale-invariant. Theorem 6.1(iii) generalizes here as:
\textbf{the ratio of consecutive block amplitudes is a
substrate-intrinsic constant}, weakly perturbed by finite-\(N\) effects.

\subsubsection{12.4 The Rank-4 Share and Finite-Rank
Truncation}\label{the-rank-4-share-and-finite-rank-truncation}

The fraction of \(\|C\|_F^2\) captured by the top-4 singular values is
approximately \(60\%\) at both scales on the real semantic substrate,
and \(41\%\) on a uniform-random control sample of \(N = 353\) unit
vectors in \(\mathbb{R}^7\). Both are well below the \(\geq 90\%\)
rank-4 dominance documented in §6 for discrete coprime-residue
substrates. The mechanism is a \textbf{finite-rank truncation of the
cosine kernel itself}, derivable from the spherical-harmonic
decomposition.

\textbf{Lemma 12.2 (Finite-rank truncation of the cosine kernel on
\(S^{d-1}\)).} \emph{Let \(K(x, y) = 1 - \langle x, y\rangle\) be the
cosine-distance kernel on the unit sphere
\(S^{d-1} \subset \mathbb{R}^d\). The spherical-harmonic decomposition
of \(K\) has nonzero content exclusively at degrees \(0\) (constant) and
\(1\) (the \(d\) directional modes), and is identically zero on all
higher degrees. Consequently, the kernel matrix
\(K_{ij} = 1 - \langle v_i, v_j\rangle\) on any sample of \(N\) unit
vectors has rank exactly \(\min(N, 1 + d)\).}

\emph{Proof.} Write \(K = J - VV^\top\) where \(J\) is the all-ones
matrix (rank 1, contributing the constant mode) and \(V\) is the
\(N \times d\) matrix of unit vectors (rank \(d\), contributing the
directional modes). \(K\) is the difference of a rank-1 and a rank-\(d\)
matrix; by sub-additivity of rank, \(\mathrm{rank}(K) \leq 1 + d\). For
\(N \geq 1 + d\) and a generic sample, equality holds. The
spherical-harmonic identity
\(1 - \langle x, y\rangle = (\text{degree-0 constant}) - \sum_{k=1}^d x_k y_k\)
decomposes into exactly \(1 + d\) rank-one summands. Higher harmonics
vanish by Gegenbauer-polynomial orthogonality applied to the linear
function \(\langle x, y\rangle\). \(\square\)

\textbf{Corollary 12.3 (Rank-4 share ceiling on \(S^{d-1}\)).} \emph{On
the cosine substrate over \(S^{d-1}\), \(C = [K, P_\tau]\) has rank at
most \(2(1 + d)\), hence the top-4 share is bounded by the largest four
squared singular values of an underlying \((1+d)\)-rank kernel
transported through the involution. For \(d = 7\),
\(\mathrm{rank}(C) \leq 16\) and the residual \(40\)--\(60\%\) of
\(\|C\|_F^2\) is distributed across the bottom \(4\) to \(12\) modes of
the kernel itself, not across an infinite overtone tower.}

\textbf{The cosine kernel does not violate \(1/k^2\) Sobolev decay; it
has no Fourier coefficients beyond degree 1 to decay at any rate.} The
triangle-wave overtone tower of Theorem 6.1 is absent on the cosine
substrate not because the regularity is degraded, but because the
kernel's spherical-harmonic support is finite. What survives the
substitution is the \emph{organization} (parity pairing, rank-4 block,
scale-invariant ratio); what does not survive is the \emph{bulk overtone
amplitude}. The \(60\%\) vs.~\(41\%\) gap between real semantic
embeddings and uniform-random controls measures the substrate's
concept-clustering structure, which concentrates additional variance
into the top modes beyond the analytical baseline.

\textbf{Observation 12.1 (Rank-4 mode persistence on continuous
Hamiltonian substrate).} The exact parity-chiral degeneracy of Theorem
6.1(i), derived for \((\mathbb{Z}/m)^\times\), holds to numerical
precision on a 768-dimensional cosine-kernel embedding of \(N = 353\)
entities under leading-eigenvector antipode involution. The rank-4 block
organization survives the substitution; the bulk overtone amplitude is
bounded by the kernel's analytical rank \(1 + d\) rather than by Sobolev
regularity. The mechanism is closed under the kink → spherical-cosine
substitution at the level of structure but not at the level of bulk
amplitude.

The empirical falsifier for this analysis is the script
\texttt{probe\_cosine\_sobolev.py} in the accompanying repository, which
verifies the spherical-harmonic rank-\((1+d)\) identity to \(10^{-13}\)
relative precision and demonstrates the kernel-induced ceiling on the
rank-4 share.

\begin{center}\rule{0.5\linewidth}{0.5pt}\end{center}

\subsection{13. Texture: An A-Optimality Scalar on the Rank-4
Subspace}\label{texture-an-a-optimality-scalar-on-the-rank-4-subspace}

Observation 12.1 establishes that the substrate's rank-4
\emph{organization} is exact and that per-mode-pair retrieval is
therefore mathematically independent across irreps. We define a scalar
that ranks a query's position within the rank-4 subspace by how well it
activates both irreps simultaneously.

\subsubsection{13.1 The Texture
Functional}\label{the-texture-functional}

Let \(\Pi_k : \mathbb{R}^4 \to \mathbb{R}^4\) denote orthogonal
projection onto the \(k\)-th mode-pair subspace (\(K = 2\) in the
present substrate). For a query vector \(q \in \mathbb{R}^4\) (the
4D-bulk projection of a 768-dimensional embedding), define the
\textbf{per-irrep tension} \[T_k(q) \;=\; \|\Pi_k\, q\|_2, \tag{13.1}\]
and the \textbf{texture}
\[\boxed{\quad \mathcal{T}(q) \;=\; \frac{K}{\sum_{k=0}^{K-1} 1/T_k(q)} \quad} \tag{13.2}\]
as the harmonic mean of per-irrep tensions.

\subsubsection{13.2 The A-Optimality
Theorem}\label{the-a-optimality-theorem}

\textbf{Theorem 13.1 (A-optimality barycenter).} \emph{Subject to
\(\sum_k T_k(q)^2 = \|q\|_2^2 = 1\), the texture \(\mathcal{T}(q)\)
attains its unique maximum at the equal-tension point
\(T_0(q) = T_1(q) = \cdots = T_{K-1}(q) = 1/\sqrt K\), with value}
\[\mathcal{T}_\text{max} \;=\; \frac{2}{\sqrt K}. \tag{13.3}\] \emph{At
\(K = 2\), \(\mathcal{T}_\text{max} = \sqrt 2 \approx 1.4142\).}

\emph{Proof.} By Cauchy-Schwarz,
\(\bigl(\sum_k T_k\bigr)\bigl(\sum_k 1/T_k\bigr) \geq K^2\), so
\(\mathcal{T}(q) = K/\sum_k(1/T_k) \leq \sum_k T_k / K\). Under
\(\sum_k T_k^2 = 1\), the Lagrangian shows \(\sum_k T_k\) is maximized
at the equal-coordinate point \(T_k = 1/\sqrt K\), giving
\(\sum_k T_k = \sqrt K\) and \(\mathcal{T}_\text{max} = 2/\sqrt K\).
Uniqueness follows from the strict concavity of the harmonic mean.
\(\square\)

\subsubsection{13.3 Why Harmonic, Not
Arithmetic}\label{why-harmonic-not-arithmetic}

The equal-tension point \(q^*\) is the \textbf{A-optimal} configuration
of the rank-4 subspace: maximal joint activation of both irreps, no
collapse onto either axis. The harmonic mean is the right aggregator
because \(1/T_k \to \infty\) as \(T_k \to 0\), which heavily penalizes
asymmetric configurations. This is the same structural reason
A-optimality is preferred over D-optimality in experimental design when
all parameters are equally important.

\textbf{Empirical saturation.} On every query tested across the
\(N = 353\) substrate, the Theorem 13.1 ceiling \(\sqrt 2\) is reached
to machine precision in \(\leq 16\) optimization steps, confirming the
bound is operational rather than asymptotic.

\begin{center}\rule{0.5\linewidth}{0.5pt}\end{center}

\subsection{14. Heisenberg-Dual
Localization}\label{heisenberg-dual-localization}

The texture barycenter \(q^*\) is the projection of a 768-dimensional
embedding; inverse-rendering requires surfacing a natural-language query
whose embedding projects to \(q^*\). The inverse map is one-to-many ---
the fibre above \(q^*\) is a dense \((768 - 4)\)-dimensional affine
subspace --- and this geometric fact has a quantitative consequence.

\subsubsection{14.1 The Obstruction}\label{the-obstruction}

\textbf{Proposition 14.1 (Fourier-dual localization).} \emph{Let
\(\mathcal{F} : L^2(\mathbb{R}^{768}) \to L^2(\mathbb{R}^4)\) denote
projection from the full embedding space to the 4D bulk of
\(C = [D_\text{sym}, P_\tau]\). For a query density \(\rho\) on
\(\mathbb{R}^{768}\) with 4D-bulk projection \(\hat\rho\) centered at
\(q^*\),}
\[\sigma_\text{topological}(\hat\rho) \cdot \sigma_\text{semantic}(\rho) \;\gtrsim\; c \tag{14.1}\]
\emph{with \(c\) a constant depending on the spectral gap of \(C\).
Delta concentration of \(\hat\rho\) at \(q^*\) implies maximally-diffuse
\(\rho\) in the semantic fibre.}

\emph{Sketch.} The projection \(\mathcal{F}\) is a linear surjection;
sharp localization in the image forces flat distribution in the fibre by
the standard Fourier-uncertainty argument applied to the fibre integral.
The constant \(c\) is set by the fibre's metric volume relative to the
spectral gap of \(C\). \(\square\)

\subsubsection{14.2 The Practical Cost}\label{the-practical-cost}

The practical consequence of (14.1) is empirically observable: two
distinct queries that reach distinct \(q^*\) points may inverse-render
to the \emph{same} natural-language sentence --- a generic trope at the
centroid of the fibre, uncorrelated with the substrate's structural
content. \textbf{Sharp localization in the topological domain forces
diffuse localization in the conjugate (semantic) domain.} \(q^*\),
though mathematically optimal, is not directly articulable: language
cannot sustain the delta.

\subsubsection{14.3 The Wave-Packet
Relaxation}\label{the-wave-packet-relaxation}

The trade-off (14.1) admits a one-parameter relaxation. Replace the
delta at \(q^*\) with a Gaussian density of width \(\sigma > 0\),
\[\hat\rho_\sigma(q) \;\propto\; \exp\!\left(-\frac{\|q - q^*\|^2}{2\sigma^2}\right). \tag{14.2}\]
\(\sigma\) interpolates between the two failure modes: \(\sigma \to 0\)
recovers the delta (sharp but linguistically collapsed);
\(\sigma \to \infty\) flattens the field into undiscriminating trope. A
sweet spot exists where the field is rich enough to support a specific
articulable rendering without losing topological specificity.

\textbf{Proposition 14.2 (Spectral-gap scaling ansatz for \(\sigma\)).}
\emph{As the substrate densifies and the spectral gap
\(\Delta = \sigma_\text{max} - \sigma_2\) grows, the basin of attraction
around each irrep axis steepens proportionally to \(\Delta\). The
Heisenberg-dual scale at which semantic coherence is preserved tightens
as}
\[\sigma(\Delta) \;=\; \sigma_\text{ref} \cdot \sqrt{\frac{\Delta_\text{ref}}{\Delta}}, \tag{14.3}\]
\emph{derived from a harmonic-oscillator analogue in which basin
curvature scales linearly with \(\Delta\) and the ground-state width
scales inversely with \(\sqrt{\text{curvature}}\).}

The single-substrate calibration at
\((\sigma_\text{ref}, \Delta_\text{ref})\) is preliminary;
multi-substrate validation across \(\geq 3\) distinct spectral states is
the next empirical check.

\begin{center}\rule{0.5\linewidth}{0.5pt}\end{center}

\subsection{15. Extrinsic Manifold
Curvature}\label{extrinsic-manifold-curvature}

Theorem 13.1 gives the mathematical ceiling
\(\mathcal{T}_\text{max} = \sqrt 2\). The wave-packet relaxation of §14
achieves an empirical rendered-texture
\(\mathcal{T}_\text{rendered} < \mathcal{T}_\text{max}\). The gap
between the two is not a bug to be gated; it is a substrate-intrinsic
geometric quantity:
\[\boxed{\quad \Delta\mathcal{T}_\text{curv}(q) \;=\; \mathcal{T}_\text{max} - \mathcal{T}_\text{rendered}(q) \;=\; \sqrt 2 - \mathcal{T}_\text{rendered}(q) \quad} \tag{15.1}\]

\subsubsection{15.1 Geometric
Interpretation}\label{geometric-interpretation}

\(\Delta\mathcal{T}_\text{curv}\) is an empirical measurement of the
\textbf{extrinsic curvature} of the semantic manifold (the
natural-language fibre) relative to the topological bulk (the 4D rank-4
subspace) at the query point. Large \(\Delta\mathcal{T}_\text{curv}\)
indicates a region where the topology prescribes a joint-irrep
activation that natural language cannot saturate --- the manifold is
highly curved in the ambient topological space at \(q^*\). Small
\(\Delta\mathcal{T}_\text{curv}\) indicates local flatness: language
nearly reaches what the math demands.

\subsubsection{15.2 Empirical Curvature
Spectrum}\label{empirical-curvature-spectrum}

On a 10-query held-out set of user intents against the \(N = 353\)
substrate, measured \(\Delta\mathcal{T}_\text{curv}\) ranges from
\(+0.30\) (flat regime: language saturates math) to \(+1.31\) (highly
curved regime: language barely reaches the prescribed joint activation).
The curvature ordering is substrate-intrinsic --- it measures the
geometry of the semantic fibre, not the quality of any individual
rendering --- and is reproducible across independent runs at fixed
\(\sigma\).

\subsubsection{15.3 The No-Gate
Discipline}\label{the-no-gate-discipline}

\textbf{Methodological discipline.} \(\Delta\mathcal{T}_\text{curv}\) is
a first-class telemetry signal to be logged, surfaced, and analyzed ---
it is not a quality gate. A fallback that substitutes the original query
when the rendered query has lower texture replaces a real measurement
with a spurious floor and blinds the instrument to the manifold's
curvature structure. The same discipline applies to the §11 Schur
compliance score and the §4 Wishart-Nyquist noise floor: all three are
measurement signals of the substrate, not pass/fail thresholds for the
user.

\subsubsection{15.4 The Stereoscopic
Coordinate}\label{the-stereoscopic-coordinate}

The triple
\((\mathcal{T}_\text{original}, \mathcal{T}_\text{rendered}, \Delta\mathcal{T}_\text{curv})\)
on a query is a \textbf{three-point coordinate on the semantic
manifold}: the original intent's position, the optimized-renderable
projection of that intent, and the extrinsic-curvature signature of the
local geometry. Together with the per-irrep tensions \((T_0, T_1)\),
this constitutes the complete first-order topological telemetry of an
intent on the rank-4 subspace.

\begin{center}\rule{0.5\linewidth}{0.5pt}\end{center}

\section{PART IV --- Synthesis}\label{part-iv-synthesis}

\subsection{16. Operator Class
Membership}\label{operator-class-membership}

The verified and candidate members of the operator class
\textbf{detailed-balance Markov generators on \(G\)-equivariant state
spaces with stratified BVP} are now the following.

{\def\LTcaptype{none} % do not increment counter
\begin{longtable}[]{@{}
  >{\raggedright\arraybackslash}p{(\linewidth - 8\tabcolsep) * \real{0.2000}}
  >{\raggedright\arraybackslash}p{(\linewidth - 8\tabcolsep) * \real{0.2000}}
  >{\raggedright\arraybackslash}p{(\linewidth - 8\tabcolsep) * \real{0.2000}}
  >{\raggedright\arraybackslash}p{(\linewidth - 8\tabcolsep) * \real{0.2000}}
  >{\raggedright\arraybackslash}p{(\linewidth - 8\tabcolsep) * \real{0.2000}}@{}}
\toprule\noalign{}
\begin{minipage}[b]{\linewidth}\raggedright
System
\end{minipage} & \begin{minipage}[b]{\linewidth}\raggedright
State space
\end{minipage} & \begin{minipage}[b]{\linewidth}\raggedright
\(G\)
\end{minipage} & \begin{minipage}[b]{\linewidth}\raggedright
Bound(s)
\end{minipage} & \begin{minipage}[b]{\linewidth}\raggedright
Status
\end{minipage} \\
\midrule\noalign{}
\endhead
\bottomrule\noalign{}
\endlastfoot
Prime gas at \(m_0 = 6\) & \(\mathbb{R}^2\) & \(\mathbb{Z}/2\) & 2 &
Verified {[}1{]}, §9 \\
Coliseum (V12 Phase A) & \([0,1]^3\) & \(S_3\) & 2, 3, 3 & Verified
§8 \\
Coprime lattices \((\mathbb{Z}/m)^\times\), \(m \geq 30\) &
\(\mathbb{R}^{\varphi(m)}\) & \((\mathbb{Z}/m)^* \times \mathbb{Z}/2\) &
rank-4 ladder & Verified §6 \\
Continuous Hamiltonian semantic substrate (\(N = 353\), 768-D) &
\(S^6 / O(7)\) & \(O(7) \times \mathbb{Z}/2\) & rank-4 (60\% share) &
Verified §12 \\
Anonymous 7-dimensional detailed-balance Markov substrate (Fancyland
LLC, internal) & \(\mathbb{R}^7\) & undeclared & parity pairing exact;
rank-4 dominance & Verified internally §9.5; substrate reserved \\
Glauber dynamics on Ising model & Lattice &
\(\Gamma \ltimes \mathbb{Z}_2^N\) & open & open \\
Metropolis-Hastings on symmetric targets & Target stabilizer & open &
open & \\
Overdamped Langevin on equivariant potentials & Potential isometry group
& open & open & \\
SYK effective Hamiltonian & \(S_N\) or subgroup & open & open & \\
Softmax attention at finite \(T\) & Attention symmetry group & open &
open & \\
\end{longtable}
}

The class is therefore \textbf{substrate-portable}: the same
character-theoretic mathematics applies to a finite-state \(S_3\)
accumulator, an arithmetic substrate of dimension \(\sim 10^5\), and a
continuous semantic embedding of dimension \(\sim 10^3\). On the third,
the same mathematics is realized as a measurement instrument.

\begin{center}\rule{0.5\linewidth}{0.5pt}\end{center}

\subsection{17. Open Questions}\label{open-questions}

Three falsifiable conjectures remain open from this work.

\textbf{Claim 6.2 (Sobolev-Kink Conjecture).} The triangle-wave block
structure holds for any \(G\)-equivariant kernel meeting the kink +
involution + circle-pullback hypotheses, abelian or non-abelian. The
natural first non-abelian test case is the \(S_3\)-equivariant Coliseum
accumulator; if Claim 6.2 holds, the same \(1/9, 9/25, 25/49\) block
ratios will appear in \([C_\text{Coliseum}, P_{(1\,2)}]\) up to
\(O(1/M)\) sampling corrections.

\textbf{Claim 7.2 (Universal coupling invariant).}
\(\kappa = \|[K, P_\tau]\|_F / \lambda_\text{Perron}(K^+) \to \sqrt{2/3}\)
on every Claim-6.2 substrate, including non-abelian \(G\). This is a
single-scalar falsifier and is much cheaper than the full
overtone-ladder protocol.

\textbf{Falsified hypothesis: Sobolev-regularity attenuation of the
rank-4 share.} A natural first hypothesis for the attenuation observed
in §12.4 is a Sobolev-regularity argument: that the rank-4 share on a
continuous substrate is bounded by
\((8/\pi^2) \cdot \|K\|_{H^{1/2}}^2 / \|K\|_{L^2}^2\), scaling with the
kernel's \(H^{1/2}\) norm. This hypothesis is falsified by
\texttt{probe\_cosine\_sobolev.py}. The cosine kernel on \(S^{d-1}\) has
no spherical-harmonic content beyond degree 1 (Lemma 12.2), so there is
no \(1/k^2\) decay to violate, and the heuristic Sobolev ratio yields
values \(> 1\) on both tent and cosine substrates --- structurally
meaningless as a ``share.'' The mechanism is instead the finite-rank
truncation of Lemma 12.2, established as a proof-grade fact in §12.4.

\begin{center}\rule{0.5\linewidth}{0.5pt}\end{center}

\subsection{18. Conclusion}\label{conclusion}

The Schur-rank bound is a-priori, testable under the Wishart-Nyquist
sampling regime, and violated only by accidental symmetry restoration
whose mechanism the FRR observation pins precisely. The Triangle-Wave
Theorem upgrades the bound to a number-theoretic spectral law with exact
rational ratios at every finite \(m\) and a single dimensionless
invariant \(\sqrt{2/3}\) derivable from two non-overlapping
character-theoretic inputs. On continuous semantic substrates the same
mathematics is realized as an extrinsic-curvature measurement instrument
--- the operator class is therefore not just a theoretical category but
a substrate-portable design pattern.

\begin{center}\rule{0.5\linewidth}{0.5pt}\end{center}

\subsection{Acknowledgments}\label{acknowledgments}

This work was carried out over a three-day window using the Lattice OS
axiomatic research protocol, in which the author proposed hypotheses and
two large language models --- Claude Opus 4 (Anthropic) and Gemini 3.1
Pro (Google DeepMind) --- served as adversarial reviewers. All claims
were resolved by computational execution, not by model assertion.

\begin{center}\rule{0.5\linewidth}{0.5pt}\end{center}

\subsection{References}\label{references}

{[}1{]} A. P. Matos, \emph{The Unity Clock: Effective Dimensional
Collapse of the Addition-Multiplication Commutator on Coprime Lattices}.
Preprint (2026).
\href{https://doi.org/10.5281/zenodo.19478727}{10.5281/zenodo.19478727}

{[}2{]} A. P. Matos, \emph{The Arithmetic Black Hole: Softmax
Thermodynamics and the Four Eigenvalue Laws of the Prime Gas}. Preprint
(2026).
\href{https://doi.org/10.5281/zenodo.19442006}{10.5281/zenodo.19442006}

{[}3{]} A. P. Matos, \emph{Active Transport on the Prime Gas: Flat-Band
Condensation, the Rabi Phase Transition, and the Arithmetic Qubit}.
Preprint (2026).
\href{https://doi.org/10.5281/zenodo.19243258}{10.5281/zenodo.19243258}

{[}4{]} A. P. Matos, \emph{Universal Two-Prime Formula for the
Coprime-Lattice Coupling Constant}. Preprint (2026).
\href{https://doi.org/10.5281/zenodo.19210625}{10.5281/zenodo.19210625}

{[}5{]} A. P. Matos, \emph{The Patient Compass: Domain Walls, Oxygen
Gating, and the 3-D Survival Manifold of 198,862 Tumors}. Preprint
(2026).
\href{https://doi.org/10.5281/zenodo.19561701}{10.5281/zenodo.19561701}

{[}6{]} A. P. Matos (Fancyland LLC), \emph{Fault-Injection-Immune
Computational Unit Using Primorial Coprime Residue Topology}. U.S.
Provisional Patent Application No.~64/031,440 (filed April 7, 2026).

{[}7{]} A. P. Matos (Fancyland LLC), \emph{Holographic Eigen-Solver
Using QM Boundary Projection on Coprime Residue Lattices}. U.S.
Provisional Patent Application No.~64/033,689 (filed April 8, 2026).

{[}8{]} A. P. Matos (Fancyland LLC), \emph{Tensegrity Interferometer:
Stereoscopic Query Resolution and Manifold-Curvature Measurement on
Continuous Hamiltonian Substrates}. U.S. Provisional Patent Application
No.~64/048,617 (filed April 24, 2026).

{[}9{]} D. Hanahan \& R. A. Weinberg, \emph{The Hallmarks of Cancer}.
Cell \textbf{100}, 57--70 (2000); updated \emph{Cell} \textbf{144},
646--674 (2011).

{[}10{]} H. Iwaniec \& E. Kowalski, \emph{Analytic Number Theory}. AMS
Colloquium Publications \textbf{53} (2004), Theorem 11.11 (Weil bound on
Kloosterman sums).

{[}11{]} A. Weil, \emph{On Some Exponential Sums}.
Proc.~Natl.~Acad.~Sci.~USA \textbf{34}, 204--207 (1948).

{[}12{]} G. H. Hardy \& E. M. Wright, \emph{An Introduction to the
Theory of Numbers} (6th ed.), Oxford University Press (2008), §5.6
(Ramanujan sums).

{[}13{]} J.-P. Serre, \emph{Linear Representations of Finite Groups}.
Graduate Texts in Mathematics \textbf{42}, Springer (1977).

{[}14{]} R. J. Muirhead, \emph{Aspects of Multivariate Statistical
Theory}. Wiley (1982), Chapters 3 and 9 (Wishart distributions).

{[}15{]} V. A. Marčenko \& L. A. Pastur, \emph{Distribution of
eigenvalues for some sets of random matrices}.
Math.~USSR-Sb.~\textbf{1}, 457--483 (1967).

{[}16{]} D. Sherrington \& S. Kirkpatrick, \emph{Solvable Model of a
Spin-Glass}. Phys.~Rev.~Lett.~\textbf{35}, 1792 (1975); A. Kitaev,
\emph{A Simple Model of Quantum Holography}. KITP Strings Seminar
(2015).

{[}17{]} A. Vaswani et al., \emph{Attention Is All You Need}. NeurIPS
(2017).

{[}18{]} J.-J. Maldacena, \emph{The Large N Limit of Superconformal
Field Theories and Supergravity}. Adv.~Theor.~Math.~Phys.~\textbf{2},
231--252 (1998).

\begin{center}\rule{0.5\linewidth}{0.5pt}\end{center}

\subsection{Appendix A:
Reproducibility}\label{appendix-a-reproducibility}

\subsubsection{A.1 Repository}\label{a.1-repository}

All math-only code reproducing the theorems of this paper is available
under MIT license at:

\begin{quote}
\textbf{\href{https://github.com/fancyland-llc/character-rank}{github.com/fancyland-llc/character-rank}}
\end{quote}

The repository contains the paper, the math-verification scripts, and no
engineering or proprietary content.

\subsubsection{A.2 Computing
Environment}\label{a.2-computing-environment}

All analyses were performed on a consumer workstation: AMD Ryzen 9 7950X
/ Windows 11. Software: Python 3.13, NumPy 2.4, SciPy 1.12. No GPU. No
external services beyond standard scientific Python.

\subsubsection{A.3 Pipeline}\label{a.3-pipeline}

The math-only reproducibility bundle exercises every theorem of this
paper via the following drivers plus one self-contained three-witness
cross-verifier.

{\def\LTcaptype{none} % do not increment counter
\begin{longtable}[]{@{}
  >{\raggedright\arraybackslash}p{(\linewidth - 2\tabcolsep) * \real{0.5000}}
  >{\raggedright\arraybackslash}p{(\linewidth - 2\tabcolsep) * \real{0.5000}}@{}}
\toprule\noalign{}
\begin{minipage}[b]{\linewidth}\raggedright
Driver
\end{minipage} & \begin{minipage}[b]{\linewidth}\raggedright
Reproduces
\end{minipage} \\
\midrule\noalign{}
\endhead
\bottomrule\noalign{}
\endlastfoot
\texttt{pytest\ sweep/} & Theorems 3.1, 5.1 --- Schur bound + entropy
decomposition (hand-built equivariant covariances) \\
\texttt{run\_holographic\_primes.py} & §6 CRT-tower spectrum, §7
\(\kappa^2 = 2/3\) scan \\
\texttt{run\_unity\_clock\_chirality.py} & §6 Lemma 6.4.1 reproduction,
§6.3 Table 6.1 fundamentals \\
\texttt{run\_horizon\_trajectory.py} & §6.3 flatline fit on 47 anchors
(3 firewall isolines) \\
\texttt{run\_overtone\_check.py} & §6.3 Table 6.1 overtone ladder (5
ratios at 8 anchors) \\
\texttt{verify\_triangle\_wave\_theorem.py} & §6 three-witness
cross-verification (Witnesses A/B/C) \\
\texttt{probe\_cosine\_sobolev.py} & §12.4 Lemma 12.2 --- cosine kernel
rank-\((1+d)\) identity, falsifier for the Sobolev-regularity hypothesis
(§17) \\
\end{longtable}
}

Total wall time end-to-end: \(\sim 20\) minutes on a Ryzen 9 7950X
single core.

\subsubsection{A.4 Three-Witness Cross-Verification of Theorem
6.1}\label{a.4-three-witness-cross-verification-of-theorem-6.1}

The verifier runs three witnesses on each anchor modulus \(m\):

\begin{itemize}
\tightlist
\item
  \textbf{Witness (A) --- Closed-form rational predictions.} Uses only
  \texttt{fractions.Fraction}. Returns
  \(\sigma_{4i}/\sigma_{4(i-1)} = (2i-1)^2/(2i+1)^2\) as exact rationals
  for the five anchor ratios. No NumPy, no SVD.
\item
  \textbf{Witness (B) --- From-scratch NumPy re-implementation.}
  Rebuilds coprime residues from \texttt{math.gcd}, the tent distance by
  direct loop, the involution by \texttt{pow(r,\ -1,\ m)}, all in stdlib
  + NumPy. Computes \([D_\text{sym}, P_\tau]\) and its top-16 singular
  values by dense \texttt{numpy.linalg.svd}.
\item
  \textbf{Witness (C) --- Production pipeline.} Calls the production
  function
  \texttt{sweep.\allowbreak{}holographic\_\allowbreak{}rank\_\allowbreak{}primes.\allowbreak{}unity\_\allowbreak{}clock\_\allowbreak{}direct\_\allowbreak{}commutator\_\allowbreak{}svd}
  --- the same function used to generate Tables 6.1 and 12.1.
\end{itemize}

{\def\LTcaptype{none} % do not increment counter
\begin{longtable}[]{@{}
  >{\raggedright\arraybackslash}p{(\linewidth - 2\tabcolsep) * \real{0.4286}}
  >{\raggedleft\arraybackslash}p{(\linewidth - 2\tabcolsep) * \real{0.5714}}@{}}
\toprule\noalign{}
\begin{minipage}[b]{\linewidth}\raggedright
Comparison
\end{minipage} & \begin{minipage}[b]{\linewidth}\raggedleft
Max absolute deviation across five ratios
\end{minipage} \\
\midrule\noalign{}
\endhead
\bottomrule\noalign{}
\endlastfoot
(B) vs (C) & \(0.000 \times 10^{0}\) --- bit-identical IEEE 754
floats \\
(B) vs external Claude separate-machine reference &
\(3.3 \times 10^{-7}\) \\
(B) vs (A) & \(3.3 \times 10^{-3}\) at
\(m = 7770, \sigma_{12}/\sigma_8\); \(\to 10^{-4}\) at \(m = 30030\)
(consistent with Theorem 6.1(iv) envelope) \\
\end{longtable}
}

The three witnesses are constructed from disjoint code paths. Witness
(A) cannot disagree with (B) or (C) modulo finite-\(\varphi\)
corrections without the rational prediction being wrong. Witness (B)
cannot disagree with (C) except through numerical error. Witness (C) is
the production code that generated every other number in §6. The
observed pattern is the theoretical one.
\textbf{\texttt{THREE-WAY\ CROSS-VERIFICATION\ PASSES.}}

\subsubsection{A.5 Scope Note on the Coliseum
Numbers}\label{a.5-scope-note-on-the-coliseum-numbers}

The Coliseum data of §8 (Table 8.1, FRR observation) was generated from
a proprietary \(N = 3\) multi-agent simulator that respects \(S_N\)
isotropy by construction. The simulator itself is \textbf{not} shipped
with this paper. What is shipped:

\begin{enumerate}
\def\labelenumi{\arabic{enumi}.}
\tightlist
\item
  The pure-math machinery of Theorems 3.1, 4.1, and 5.1, validated on
  hand-built \(G\)-equivariant covariances with analytically-known
  spectra (\texttt{sweep/test\_character\_rank.py},
  \texttt{sweep/test\_holographic\_rank.py}).
\item
  The prime-gas Dirichlet-character machinery
  (\texttt{sweep/holographic\_rank\_primes.py}), arithmetic on its own
  and not requiring any simulator.
\item
  Drivers reproducing §6, §7, and §12.1--12.3 directly from closed-form
  formulas.
\end{enumerate}

Independent re-derivation of §8's specific numbers requires a
third-party generator preserving \(G\)-equivariance; the tests here pin
the math any such generator's output must obey.

\subsubsection{A.6 Scope Note on the Continuous Substrate
Numbers}\label{a.6-scope-note-on-the-continuous-substrate-numbers}

The continuous Hamiltonian substrate of §§10--15 (Observation 12.1,
Table 12.1, the §15 curvature spectrum) was constructed on a proprietary
768-dimensional semantic embedding with a learned 7-feature compression.
The embedding model and compression are \textbf{not} shipped, and the
texture optimizer, Gaussian wave-packet relaxation, and composed
measurement pipeline are intentionally omitted from this preprint per
Patent Notice {[}8{]}. What is reproducible from the math-only bundle:

\begin{enumerate}
\def\labelenumi{\arabic{enumi}.}
\tightlist
\item
  Theorem 13.1 (A-optimality of the harmonic mean) is provable from
  Cauchy-Schwarz and reproduced in \texttt{verify\_texture\_aopt.py} on
  hand-built rank-4 subspaces with analytically-known maxima.
\item
  Proposition 14.1 (Heisenberg-dual obstruction) is provable from
  standard Fourier uncertainty applied to the fibre integral; the
  constant \(c\) depends on the spectral gap of the substrate.
\item
  The §15 curvature spectrum is substrate-specific empirical data; its
  qualitative shape (substrate-intrinsic, reproducible across runs at
  fixed \(\sigma\)) is what the math predicts.
\end{enumerate}

Apparatus practicing the §13--§15 measurement program in the manner
claimed in {[}8{]} requires a license from Fancyland LLC.

\subsubsection{A.7 Relationship to Companion
Patents}\label{a.7-relationship-to-companion-patents}

The §6, §7, and §12 results were originally identified using the
proprietary spectral methods of {[}6{]}, {[}7{]}, and {[}8{]}. Those
methods accelerated discovery but are not required for verification: all
theorems and quantitative claims of this paper are independently
reproducible using only standard representation theory, NumPy linear
algebra, and the public math-only bundle at
\href{https://github.com/fancyland-llc/character-rank}{github.com/fancyland-llc/character-rank}.

\vspace{3em}
\begin{center}
\rule{0.4\textwidth}{0.5pt}

\vspace{0.8em}

\textit{Fancyland LLC --- Lattice OS research infrastructure.}

\vspace{0.3em}

\textit{The rabbit has been caught.}
\end{center}

\end{document}
